\chapter{Ruang Metrik} \label{ms:chapter}

%%%%%%%%%%%%%%%%%%%%%%%%%%%%%%%%%%%%%%%%%%%%%%%%%%%%%%%%%%%%%%%%%%%%%%%%%%%%%%

\section{Ruang metrik}
\label{sec:metric}

\sectionnotes{1,5 kuliah}

Seperti disebutkan dalam pendahuluan, gagasan utama dalam analisis adalah
menentukan limit.  Dalam \chapterref{seq:chapter}, kita mempelajari limit
barisan bilangan real.  Dalam \chapterref{lim:chapter}, kita mempelajari limit
fungsi ketika suatu variabel real mendekati bilangan real tertentu.

Kita ingin menentukan limit dalam konteks yang lebih rumit.  Sebagai
contoh, kita ingin membahas barisan titik dalam ruang berdimensi 3.
Kita ingin mendefinisikan fungsi kontinu dengan beberapa variabel.
Kita bahkan ingin mendefinisikan fungsi pada ruang yang agak lebih sulit
dideskripsikan, seperti permukaan bumi.  Di sana pun kita tetap ingin
membahas limit.

Terakhir, dalam \chapterref{fs:chapter} kita telah melihat limit suatu
barisan fungsi.
Kita ingin menyatukan semua pengertian ini agar tidak perlu
membuktikan ulang teorema demi teorema dalam setiap konteks.  Konsep
ruang metrik merupakan alat dasar yang ampuh dalam analisis.
Meskipun tidak memadai untuk mendeskripsikan setiap jenis limit yang
dijumpai dalam analisis modern, konsep ini membawa kita sangat jauh.

\begin{defn}
Misalkan $X$ suatu himpunan, dan misalkan
$d \colon X \times X \to \R$
suatu fungsi sedemikian sehingga untuk semua $x,y,z \in X$ berlaku
\begin{enumerate}[(i)]
%
\item \label{metric:pos}
%mbxSTARTIGNORE
\makebox[2.8in][l]{$d(x,y) \geq 0$.}
%mbxENDIGNORE
%mbxlatex $d(x,y) \geq 0$. \qquad
(\emph{sifat taknegatif})\index{sifat taknegatif metrik}
%
\item \label{metric:zero}
%mbxSTARTIGNORE
\makebox[2.8in][l]{$d(x,y) = 0$ jika dan hanya jika $x = y$.}
%mbxENDIGNORE
%mbxlatex $d(x,y) = 0$ jika dan hanya jika $x = y$. \qquad
(\emph{\myindex{identitas unsur yang tak terbedakan}})
%
\item \label{metric:com}
%mbxSTARTIGNORE
\makebox[2.8in][l]{$d(x,y) = d(y,x)$.} 
%mbxENDIGNORE
%mbxlatex $d(x,y) = d(y,x)$. \qquad
(\emph{simetri})\index{simetri metrik}
%
\item \label{metric:triang}
%mbxSTARTIGNORE
\makebox[2.8in][l]{$d(x,z) \leq d(x,y)+ d(y,z)$.}
%mbxENDIGNORE
%mbxlatex $d(x,z) \leq d(x,y)+ d(y,z)$. \qquad
(\emph{\myindex{ketaksamaan segitiga}})
\end{enumerate}
Pasangan $(X,d)$ disebut \emph{\myindex{ruang metrik}}.  Fungsi $d$
disebut \emph{\myindex{metrik}} atau \emph{\myindex{fungsi jarak}}.
Jika metriknya jelas dari konteks, terkadang kita cukup menulis $X$
sebagai ruang metrik alih-alih $(X,d)$.
\end{defn}

Secara geometris, $d$ adalah jarak antara dua titik.
Butir \ref{metric:pos}--\ref{metric:com} memiliki tafsiran geometris yang
jelas: Jarak selalu taknegatif, satu-satunya titik yang berjarak 0 dari
$x$ adalah $x$ itu sendiri, dan akhirnya jarak dari $x$ ke $y$ sama dengan
jarak dari $y$ ke $x$.  Ketaksamaan segitiga \ref{metric:triang}
memiliki tafsiran yang diberikan dalam
\figureref{fig:mstriang}.
\begin{myfigureht}
\myincludepdft{ms-triang}{%
Segitiga dengan tiga titik sudut x, y, dan z.
Jarak antara x dan z ditandai sebagai
d dari x koma z, jarak dari x ke y
ditandai sebagai d dari x koma y, dan jarak
dari y ke z ditandai sebagai d dari y koma z.
Panah dari x ke z berlabel ``lebih pendek'', sedangkan
panah yang bergerak dari x ke y lalu berbelok ke z
berlabel ``lebih panjang''.}
\caption{Diagram ketaksamaan segitiga dalam ruang metrik.\label{fig:mstriang}}
\end{myfigureht}

Untuk keperluan ilustrasi, akan lebih mudah jika gambar dan diagram dibuat
pada bidang dengan jarak Euklides sebagai metrik.  Namun, itu hanyalah
satu contoh khusus ruang metrik.  Fakta bahwa suatu pernyataan tampak
jelas dari sebuah gambar tidak berarti bahwa pernyataan itu benar dalam
setiap ruang metrik.  Intuisi dari geometri Euklides dapat menyesatkan
kita, padahal konsep ruang metrik jauh lebih umum.

Berikut beberapa contoh ruang metrik.

\begin{example}
Himpunan bilangan real $\R$ merupakan ruang metrik dengan metrik
\begin{equation*}
d(x,y) \coloneqq \sabs{x-y} .
\end{equation*}
Butir \ref{metric:pos}--\ref{metric:com} dalam definisi tersebut
mudah diperiksa.  Ketaksamaan segitiga \ref{metric:triang} langsung
mengikuti ketaksamaan segitiga biasa untuk bilangan real:
\begin{equation*}
d(x,z) = \sabs{x-z} = 
\sabs{x-y+y-z} \leq
\sabs{x-y}+\sabs{y-z} =
d(x,y)+ d(y,z) .
\end{equation*}
Metrik ini adalah \emph{\myindex{metrik standar pada $\R$}}.  Jika kita
membahas $\R$ sebagai ruang metrik tanpa menyebut metrik tertentu,
yang dimaksud adalah metrik ini.
\end{example}

\begin{example}
Kita juga dapat mendefinisikan metrik lain pada himpunan bilangan real.
Sebagai contoh, ambil himpunan bilangan real $\R$ bersama metrik
\begin{equation*}
d(x,y) \coloneqq
\frac{\sabs{x-y}}{\sabs{x-y}+1} .
\end{equation*}
Butir \ref{metric:pos}--\ref{metric:com} sekali lagi mudah diperiksa.
Ketaksamaan segitiga \ref{metric:triang} sedikit lebih sulit.
Perhatikan bahwa $d(x,y) = \varphi\bigl(\sabs{x-y}\bigr)$, dengan
$\varphi(t) =
\frac{t}{t+1}$, dan bahwa $\varphi$ merupakan fungsi naik
(turunannya positif; lihat \figureref{fig:tovertp1}).
Oleh karena itu,
\begin{equation*}
\begin{split}
d(x,z) = \varphi\bigl(\sabs{x-z}\bigr)
& = 
\varphi\bigl(\sabs{x-y+y-z}\bigr) \\
& \leq
\varphi\bigl(\sabs{x-y}+\sabs{y-z}\bigr)
\\
& =
\frac{\sabs{x-y}+\sabs{y-z}}{\sabs{x-y}+\sabs{y-z}+1} \\
& =
\frac{\sabs{x-y}}{\sabs{x-y}+\sabs{y-z}+1} +
\frac{\sabs{y-z}}{\sabs{x-y}+\sabs{y-z}+1}
\\
& \leq
\frac{\sabs{x-y}}{\sabs{x-y}+1} +
\frac{\sabs{y-z}}{\sabs{y-z}+1} =
d(x,y)+ d(y,z) .
\end{split}
\end{equation*}
Dengan demikian, fungsi $d$ merupakan metrik dan memberikan
contoh metrik nonstandar pada $\R$.  Dengan metrik ini,
$d(x,y) < 1$ untuk semua $x,y \in \R$.  Artinya,
setiap pasangan titik berjarak kurang dari 1 satuan.
\begin{myfigureht}
\myincludegraphics{tovertp1graph}{%
Grafik suatu fungsi yang mula-mula menanjak secara diagonal
dari sumbu horizontal, lalu perlahan melandai
hingga hampir horizontal di ujung kanan grafik;
sebuah asimtot horizontal di atas grafik
ditandai dengan garis putus-putus.}
\caption{Grafik $\frac{t}{t+1}$ untuk $t$ positif dengan asimtot pada 1.\label{fig:tovertp1}}
\end{myfigureht}
\end{example}

Suatu ruang metrik yang penting adalah
\emph{\myindex{ruang Euklides}} berdimensi $n$,
\glsadd{not:euclidspace}
$\R^n = \R \times \R \times \cdots \times \R$.   Kita menggunakan notasi berikut
untuk titik: $x =(x_1,x_2,\ldots,x_n) \in \R^n$.  Kita tidak menuliskan
$\vec{x}$ ataupun $\mathbf{x}$ untuk sebuah titik di $\R^n$ sebagaimana lazim dalam
kalkulus multivariabel;
kita cukup memberinya nama seperti $x$ dan mengingat bahwa $x$
merupakan unsur $\R^n$.
Kita juga
cukup menuliskan $0 \in \R^n$ untuk menyatakan titik $(0,0,\ldots,0)$.  Sebelum
menjadikan $\R^n$ sebuah ruang metrik, kita membuktikan suatu ketaksamaan penting, yaitu
ketaksamaan Cauchy--Schwarz.

\begin{lemma}[\myindex{ketaksamaan Cauchy--Schwarz}%
\footnote{%
Kadang-kadang ketaksamaan ini disebut \myindex{ketaksamaan Cauchy--Bunyakovsky--Schwarz}.
\href{https://en.wikipedia.org/wiki/Hermann_Schwarz}{Karl Hermann Amandus
Schwarz} (1843--1921) adalah matematikawan Jerman dan
\href{https://en.wikipedia.org/wiki/Viktor_Bunyakovsky}{Viktor Yakovlevich
Bunyakovsky} (1804--1889) adalah matematikawan Ukraina.  Ketaksamaan yang
kita nyatakan sebenarnya semestinya disebut ketaksamaan Cauchy, karena
Bunyakovsky dan Schwarz memberikan bukti untuk versi berdimensi takhingga.}]
Misalkan $x =(x_1,x_2,\ldots,x_n) \in \R^n$, $y =(y_1,y_2,\ldots,y_n) \in \R^n$.
Maka
\begin{equation*}
{\biggl( \sum_{k=1}^n x_k y_k \biggr)}^2
\leq
\biggl(\sum_{k=1}^n x_k^2 \biggr)
\biggl(\sum_{k=1}^n y_k^2 \biggr) .
\end{equation*}
\end{lemma}

\begin{proof}
Kuadrat bilangan real selalu taknegatif.  Oleh karena itu, jumlah kuadrat
juga taknegatif:
\begin{equation*}
\begin{split}
0 & \leq 
\sum_{k=1}^n \sum_{\ell=1}^n {(x_k y_\ell - x_\ell y_k)}^2
\\
& =
\sum_{k=1}^n \sum_{\ell=1}^n \bigl( x_k^2 y_\ell^2 + x_\ell^2 y_k^2 - 2 x_k
x_\ell y_k
y_\ell \bigr)
\\
& =
\biggl( \sum_{k=1}^n x_k^2 \biggr)
\biggl( \sum_{\ell=1}^n y_\ell^2 \biggr)
+
\biggl( \sum_{k=1}^n y_k^2 \biggr)
\biggl( \sum_{\ell=1}^n x_\ell^2 \biggr)
-
2
\biggl( \sum_{k=1}^n x_k y_k \biggr)
\biggl( \sum_{\ell=1}^n x_\ell y_\ell \biggr) .
\end{split}
\end{equation*}
Dengan mengganti nama indeks dan membagi dengan 2, kita memperoleh tepat ketaksamaan yang diinginkan,
\begin{equation*}
0 \leq 
\biggl( \sum_{k=1}^n x_k^2 \biggr)
\biggl( \sum_{k=1}^n y_k^2 \biggr)
-
{\biggl( \sum_{k=1}^n x_k y_k \biggr)}^2 .
\qedhere
\end{equation*}
\end{proof}

\begin{example}
Kita akan membangun
metrik standar\index{metrik standar pada $\R^n$} untuk $\R^n$.  Definisikan
\begin{equation*}
d(x,y) \coloneqq
\sqrt{
{(x_1-y_1)}^2 + 
{(x_2-y_2)}^2 + 
\cdots +
{(x_n-y_n)}^2
} =
\sqrt{
\sum_{k=1}^n
{(x_k-y_k)}^2 
} .
\end{equation*}
Untuk $n=1$, yakni garis real, metrik ini sama dengan metrik yang didefinisikan di atas.
Untuk $n > 1$,
seperti sebelumnya, satu-satunya bagian definisi yang agak rumit untuk diperiksa adalah ketaksamaan segitiga.
Perhitungannya lebih sederhana jika kita bekerja dengan kuadrat metrik.  Dalam
penaksiran berikut, perhatikan penggunaan ketaksamaan Cauchy--Schwarz.
\begin{equation*}
\begin{split}
{\bigl(d(x,z)\bigr)}^2 & =
\sum_{k=1}^n
{(x_k-z_k)}^2 
\\
& =
\sum_{k=1}^n
{(x_k-y_k+y_k-z_k)}^2 
\\
& =
\sum_{k=1}^n
\Bigl(
{(x_k-y_k)}^2+{(y_k-z_k)}^2 + 2(x_k-y_k)(y_k-z_k)
\Bigr)
\\
& =
\sum_{k=1}^n
{(x_k-y_k)}^2
+
\sum_{k=1}^n
{(y_k-z_k)}^2 
+
2
\sum_{k=1}^n
(x_k-y_k)(y_k-z_k)
\\
& \leq
\sum_{k=1}^n
{(x_k-y_k)}^2
+
\sum_{k=1}^n
{(y_k-z_k)}^2 
+
2
\sqrt{
\sum_{k=1}^n
{(x_k-y_k)}^2
\sum_{k=1}^n
{(y_k-z_k)}^2
}
\\
& =
{\left(
\sqrt{
\sum_{k=1}^n
{(x_k-y_k)}^2
}
+
\sqrt{
\sum_{k=1}^n
{(y_k-z_k)}^2 
}
\right)}^2
=
{\bigl( d(x,y) + d(y,z) \bigr)}^2 .
\end{split}
\end{equation*}
Karena akar kuadrat merupakan fungsi naik,
ketaksamaan tetap berlaku ketika kita 
mengambil akar kuadrat kedua ruas, sehingga diperoleh ketaksamaan segitiga.
\end{example}

\begin{example}
Himpunan bilangan kompleks $\C$ adalah himpunan bilangan $z = x+iy$, dengan $x$
dan $y$ berada di $\R$.  Dengan menetapkan $i^2 = -1$, kita menjadikan $\C$
sebuah medan.
Untuk keperluan menentukan limit,
himpunan $\C$ dipandang sebagai ruang metrik $\R^2$, dengan $z=x+iy \in \C$
bersesuaian dengan $(x,y) \in \R^2$.
Untuk $z=x+iy$, definisikan \emph{\myindex{modulus kompleks}}\index{modulus}
sebagai $\sabs{z} \coloneqq \sqrt{x^2+y^2}$.
Maka untuk dua bilangan kompleks
$z_1 = x_1 + iy_1$ dan $z_2 = x_2 + iy_2$, jaraknya adalah
\begin{equation*}
d(z_1,z_2) = \sqrt{{(x_1-x_2)}^2+ {(y_1-y_2)}^2} = \sabs{z_1-z_2}.
\end{equation*}
Selain itu, ketika bekerja dengan bilangan kompleks,
sering kali lebih mudah menuliskan metrik dalam bentuk
\emph{\myindex{konjugat kompleks}}: Konjugat dari $z=x+iy$
adalah $\bar{z} \coloneqq x-iy$.  Dengan demikian,
${\sabs{z}}^2 = x^2 +y^2 = z\bar{z}$, sehingga ${\sabs{z_1-z_2}}^2 =
(z_1-z_2)\overline{(z_1-z_2)}$.
\end{example}

\begin{example}
Contoh yang perlu selalu diingat adalah
\emph{\myindex{metrik diskret}}.
Untuk sebarang himpunan $X$, definisikan
\begin{equation*}
d(x,y) \coloneqq
\begin{cases}
1 & \text{jika } x \neq y, \\
0 & \text{jika } x = y.
\end{cases}
\end{equation*}
Dengan kata lain, jarak di antara setiap dua titik yang berbeda adalah sama.  Jika $X$
merupakan himpunan berhingga, kita dapat membuat diagram; lihat, misalnya,
\figureref{fig:msdiscmetric}.  Tentu saja, jarak dalam diagram tersebut
bukanlah jarak Euklides biasa pada bidang.
Keadaannya menjadi lebih rumit ketika $X$ merupakan himpunan takhingga,
seperti himpunan bilangan real.
\begin{myfigureht}
\myincludepdft{msdiscmetric}{%
Diagram lima titik a, b, c, d, dan e.  Setiap titik
dihubungkan dengan setiap titik lainnya oleh sebuah garis, dan setiap garis
diberi label angka 1.}
\caption{Contoh ruang metrik diskret $\{ a,b,c,d,e \}$; jarak
di antara setiap dua titik yang berbeda adalah $1$.\label{fig:msdiscmetric}}
\end{myfigureht}

Meskipun contoh khusus ini mungkin jarang muncul dalam praktik, contoh ini
memberikan \myquote{uji kewajaran} yang berguna.  Jika kita mengajukan suatu
pernyataan tentang ruang metrik, cobalah mengujinya dengan metrik diskret.
Pembuktian bahwa $(X,d)$ memang merupakan ruang metrik diserahkan sebagai latihan.
\end{example}

\begin{example} \label{example:msC01}
Misalkan $C\bigl([a,b],\R\bigr)$\glsadd{not:contfuncs} adalah himpunan
fungsi kontinu bernilai real pada interval $[a,b]$.
Definisikan metrik pada $C\bigl([a,b],\R\bigr)$ dengan
\begin{equation*}
d(f,g) \coloneqq \sup_{x \in [a,b]} \babs{f(x)-g(x)} .
\end{equation*}
Kita periksa sifat-sifatnya.  Pertama, $d(f,g)$ berhingga karena
$\babs{f(x)-g(x)}$ merupakan fungsi kontinu pada interval tertutup dan terbatas
$[a,b]$, sehingga fungsi tersebut terbatas.
Jelas bahwa $d(f,g) \geq 0$;
nilai itu merupakan supremum bilangan-bilangan taknegatif.  Jika $f = g$,
maka $\babs{f(x)-g(x)} = 0$ untuk semua $x$, sehingga $d(f,g) = 0$.  Sebaliknya,
jika $d(f,g) = 0$, maka untuk setiap $x$ berlaku $\babs{f(x)-g(x)} \leq d(f,g) = 0$,
sehingga $f(x) = g(x)$ untuk semua $x$, dan karenanya $f=g$.  Kesamaan
$d(f,g) = d(g,f)$ juga sama jelasnya.  Untuk membuktikan ketaksamaan segitiga,
kita menggunakan ketaksamaan segitiga biasa.
\begin{equation*}
\begin{split}
d(f,g) & =
\sup_{x \in [a,b]} \babs{f(x)-g(x)} =
\sup_{x \in [a,b]} \babs{f(x)-h(x)+h(x)-g(x)}
\\
& \leq
\sup_{x \in [a,b]} \Bigl( \babs{f(x)-h(x)}+\babs{h(x)-g(x)} \Bigr)
\\
& \leq
\sup_{x \in [a,b]} \babs{f(x)-h(x)}+
\sup_{x \in [a,b]} \babs{h(x)-g(x)} = d(f,h) + d(h,g) .
\end{split}
\end{equation*}
Ketika $C\bigl([a,b],\R\bigr)$ dipandang sebagai ruang metrik tanpa menyebutkan metriknya,
yang dimaksud adalah metrik khusus ini.
Perhatikan bahwa $d(f,g) = \snorm{f-g}_{[a,b]}$, yaitu norma seragam yang didefinisikan dalam \defnref{def:unifnorm}.

Contoh ini mungkin tampak abstrak pada awalnya, tetapi ternyata bekerja dengan
ruang seperti $C\bigl([a,b],\R\bigr)$ benar-benar merupakan inti dari sebagian besar
analisis modern.  Dengan memandang himpunan fungsi sebagai ruang metrik, kita dapat
mengabstraksikan banyak rincian teknis yang merepotkan dan membuktikan hasil-hasil kuat seperti
\hyperref[thm:fs:picard]{teorema Picard} dengan lebih sedikit usaha.
\end{example}

\begin{example}
Contoh ruang metrik berguna lainnya adalah permukaan bola
dengan metrik yang biasanya disebut \emph{\myindex{jarak lingkaran besar}}.
Misalkan $S^2$ adalah \myindex{permukaan bola satuan}\index{permukaan bola} di $\R^3$,
yaitu $S^2 \coloneqq \{ x \in \R^3 : x_1^2+x_2^2+x_3^2 = 1 \}$.
Ambil $x$ dan $y$ di $S^2$, lalu pandang vektor dari titik asal menuju $x$
dan vektor dari titik asal menuju $y$.
Misalkan $\theta$ adalah sudut di antara kedua vektor tersebut.
Kemudian definisikan $d(x,y) \coloneqq \theta$.  Lihat \figureref{fig:spheremetric}.
Hukum kosinus dari kalkulus vektor memberikan
$d(x,y) = \arccos(x_1y_1+x_2y_2+x_3y_3)$.
Relatif mudah untuk melihat bahwa fungsi ini memenuhi tiga sifat pertama
suatu metrik.
Ketaksamaan segitiga lebih sulit dibuktikan dan memerlukan sedikit lebih banyak
trigonometri serta aljabar linear daripada yang ingin kita gunakan saat ini,
jadi kita biarkan tanpa bukti.

\begin{myfigureht}
\myincludepdft{spheremetric}{%
Diagram permukaan bola dalam ruang berdimensi 3
yang ditandai sebagai S pangkat 2.  Titik asal ditandai dengan 0,
dan terdapat dua titik pada permukaan bola, x dan y.  Terdapat
garis dari titik asal ke x, garis dari titik asal ke y,
dan sudut di antara kedua garis tersebut diberi label theta.}
\caption{Jarak lingkaran besar pada permukaan
bola satuan.\label{fig:spheremetric}}
\end{myfigureht}

Jarak ini merupakan jarak terpendek di antara titik-titik pada permukaan bola jika
kita hanya diperbolehkan bergerak di permukaannya.  Hal ini mudah
digeneralisasi untuk jari-jari sebarang.  Jika kita mengambil permukaan bola berjari-jari
$r$, jaraknya kita tetapkan sebagai $d(x,y) \coloneqq r \theta$.  Sebagai contoh, inilah
cara baku untuk mengukur jarak pada
permukaan bumi, misalnya jarak yang ditempuh pesawat dari London ke
Los Angeles.
\end{example}

Sering kali berguna untuk memandang suatu himpunan bagian dari ruang metrik yang lebih besar
sebagai ruang metrik tersendiri.  Kita memperoleh proposisi berikut, yang
memiliki bukti langsung.

\begin{prop}
Misalkan $(X,d)$ adalah ruang metrik dan $Y \subset X$.  Maka pembatasan
$d|_{Y \times Y}$ merupakan metrik pada $Y$.
\end{prop}

\begin{defn}
Jika $(X,d)$ merupakan ruang metrik, $Y \subset X$, dan $d' \coloneqq d|_{Y \times Y}$,
maka $(Y,d')$ disebut \emph{\myindex{subruang}} dari $(X,d)$.
\end{defn}

Lazimnya kita cukup menulis $d$ untuk metrik pada $Y$, sebab metrik tersebut
merupakan pembatasan metrik pada $X$.  Terkadang kita mengatakan bahwa $d'$
adalah \emph{\myindex{metrik subruang}} dan $Y$ memiliki
\emph{\myindex{topologi subruang}}.

\medskip

Suatu himpunan bagian dari bilangan real
bersifat terbatas apabila semua unsurnya berjarak paling besar suatu nilai tetap
dari 0.
Dalam ruang metrik sebarang, mungkin tidak ada titik tetap alami 0,
tetapi hal ini tidak menjadi masalah untuk konsep keterbatasan.

\begin{defn}
Misalkan $(X,d)$ adalah ruang metrik.  Suatu himpunan bagian $S \subset X$ disebut
\emph{terbatas}\index{himpunan terbatas} jika terdapat $p \in X$ dan
$B \in \R$ sedemikian sehingga
\begin{equation*}
d(p,x) \leq B \quad \text{untuk semua } x \in S.
\end{equation*}
Kita mengatakan bahwa $(X,d)$ terbatas jika $X$ sendiri merupakan himpunan bagian terbatas.
\end{defn}

Sebagai contoh, himpunan bilangan real dengan metrik standar bukanlah
ruang metrik terbatas.  Tidak sulit untuk melihat bahwa suatu
himpunan bagian bilangan real bersifat terbatas dalam
pengertian \chapterref{rn:chapter} jika dan hanya jika himpunan itu terbatas sebagai himpunan bagian dari
ruang metrik bilangan real dengan metrik standar.

Di sisi lain, jika kita mengambil bilangan real dengan metrik diskret,
kita memperoleh ruang metrik terbatas.  Bahkan, sebarang himpunan dengan
metrik diskret bersifat terbatas.

Ada beberapa cara ekuivalen lain untuk memperumum konsep keterbatasan,
yang diserahkan sebagai latihan.  Anggap $X$ tidak kosong untuk menghindari persoalan teknis.
Maka, untuk himpunan bagian tidak kosong $S \subset X$, sifat terbatas ekuivalen dengan setiap kondisi berikut:
\begin{enumerate}[(i)]
\item
Untuk setiap $p \in X$, terdapat $B > 0$ sedemikian sehingga $d(p,x) \leq B$ untuk
semua $x \in S$.
\item
\glsadd{not:diam}%
$\operatorname{diam}(S) \coloneqq \sup \bigl\{ d(x,y) : x,y \in S \bigr\} < \infty$.
\end{enumerate}
Besaran $\operatorname{diam}(S)$ disebut
\emph{\myindex{diameter}} suatu himpunan dan biasanya hanya didefinisikan untuk
himpunan yang tidak kosong.

\subsection{Latihan}

\begin{exercise}
Tunjukkan bahwa untuk setiap himpunan $X$, metrik diskret ($d(x,y) = 1$ jika $x \neq y$ dan
$d(x,x) = 0$) memang menjadikan $(X,d)$ suatu ruang metrik.
\end{exercise}

\begin{exercise}
Misalkan $X \coloneqq \{ 0 \}$ suatu himpunan.  Dapatkah Anda menjadikannya ruang metrik?
\end{exercise}

\begin{exercise}
Misalkan $X \coloneqq \{ a, b \}$ suatu himpunan.  Dapatkah Anda menjadikannya dua ruang
metrik yang berbeda?  (definisikan dua metrik yang berbeda padanya)
\end{exercise}

\begin{exercise}
Misalkan himpunan $X \coloneqq \{ A, B, C \}$ mewakili 3 gedung di kampus.  Misalkan kita
ingin jarak menyatakan waktu yang diperlukan untuk berjalan kaki dari satu gedung ke
gedung lainnya.
Perjalanan antara gedung $A$ dan $B$ memerlukan 5 menit ke arah mana pun.  Akan tetapi,
gedung $C$ berada di atas bukit dan diperlukan 10 menit dari $A$ serta 15 menit
dari $B$ untuk mencapai $C$.  Sebaliknya, diperlukan 5 menit untuk berjalan
dari $C$ ke $A$ dan 7 menit dari $C$ ke $B$, karena kita berjalan
menuruni bukit.  Apakah jarak-jarak ini mendefinisikan metrik?  Jika ya, buktikan; jika tidak, jelaskan
alasannya.
\end{exercise}

\begin{exercise}
Misalkan $(X,d)$ suatu ruang metrik dan
$\varphi \colon [0,\infty) \to \R$ suatu
fungsi naik sedemikian sehingga
$\varphi(t) \geq 0$ untuk setiap $t$ dan $\varphi(t) = 0$ jika dan hanya jika
$t=0$.  Andaikan pula $\varphi$ bersifat \emph{\myindex{subaditif}},
yaitu $\varphi(s+t) \leq \varphi(s)+\varphi(t)$.
Tunjukkan bahwa dengan $d'(x,y) \coloneqq \varphi\bigl(d(x,y)\bigr)$, kita memperoleh
ruang metrik baru $(X,d')$.
\end{exercise}

\begin{exercise} \label{exercise:mscross}
Misalkan $(X,d_X)$ dan $(Y,d_Y)$ ruang-ruang metrik.
\begin{enumerate}[a)]
\item
Tunjukkan bahwa $(X \times Y,d)$ dengan
$d\bigl( (x_1,y_1), (x_2,y_2) \bigr) \coloneqq d_X(x_1,x_2) + d_Y(y_1,y_2)$ merupakan
ruang metrik.
\item
Tunjukkan bahwa $(X \times Y,d)$ dengan
$d\bigl( (x_1,y_1), (x_2,y_2) \bigr) \coloneqq \max \bigl\{ d_X(x_1,x_2) ,
d_Y(y_1,y_2) \bigr\}$ merupakan
ruang metrik.
\end{enumerate}
\end{exercise}

\begin{exercise}
Misalkan $X$ himpunan fungsi-fungsi kontinu pada $[0,1]$.  Misalkan $\varphi \colon
[0,1] \to (0,\infty)$ kontinu.  Definisikan
\begin{equation*}
d(f,g) \coloneqq \int_0^1 \babs{f(x)-g(x)}\varphi(x)\,dx .
\end{equation*}
Tunjukkan bahwa $(X,d)$ merupakan ruang metrik.
\end{exercise}

\begin{exercise} \label{exercise:mshausdorffpseudo}
Misalkan $(X,d)$ suatu ruang metrik.  Untuk himpunan bagian tak kosong dan terbatas $A$ dan $B$, definisikan
\begin{equation*}
d(x,B) \coloneqq \inf \bigl\{ d(x,b) : b \in B \bigr\}
\qquad \text{dan} \qquad
d(A,B) \coloneqq \sup \bigl\{ d(a,B) : a \in A \bigr\} .
\end{equation*}
Sekarang definisikan \emph{\myindex{metrik Hausdorff}} sebagai
\begin{equation*}
d_H(A,B) \coloneqq \max \bigl\{ d(A,B) , d(B,A) \bigr\} .
\end{equation*}
Catatan: $d_H$ dapat didefinisikan untuk sebarang himpunan bagian tak kosong jika kita mengizinkan
bilangan real diperluas.
\begin{enumerate}[a)]
\item
Misalkan $Y \subset \sP(X)$ himpunan semua himpunan bagian tak kosong dan terbatas.
Buktikan bahwa
$(Y,d_H)$ merupakan apa yang kita sebut \emph{\myindex{ruang pseudometrik}}:
$d_H$ memenuhi sifat-sifat metrik
\ref{metric:pos},
\ref{metric:com}, 
\ref{metric:triang}, dan juga
$d_H(A,A) = 0$ untuk setiap $A \in Y$.
\item
Tunjukkan melalui contoh bahwa $d$ sendiri tidak simetris,
yaitu $d(A,B) \neq d(B,A)$.
\item
Carilah suatu ruang metrik $X$ dan dua himpunan bagian tak kosong
dan terbatas yang berbeda, $A$ dan $B$, sedemikian sehingga $d_H(A,B) = 0$.
\end{enumerate}
\end{exercise}

\begin{exercise}
Misalkan $(X,d)$ ruang metrik tak kosong dan $S \subset X$ suatu himpunan bagian.  Buktikan:
\begin{enumerate}[a)]
\item
$S$ terbatas jika dan hanya jika
untuk setiap $p \in X$, terdapat $B > 0$ sedemikian sehingga $d(p,x) \leq B$ untuk
setiap $x \in S$.
\item
Himpunan $S$ yang tak kosong terbatas jika dan hanya jika
$\operatorname{diam}(S) \coloneqq \sup \{ d(x,y) : x,y \in S \} < \infty$.
\end{enumerate}
\end{exercise}

\begin{exercise}
\pagebreak[2]
\leavevmode
\begin{enumerate}[a)]
\item
Dengan bekerja di $\R$, hitunglah $\operatorname{diam}\bigl([a,b]\bigr)$.
\item
Dengan bekerja di $\R^n$, untuk setiap $r > 0$, misalkan $B_r \coloneqq \{ x \in \R^n : x_1^2+x_2^2+\cdots+x_n^2
< r^2 \}$.  Hitunglah $\operatorname{diam}(B_r)$.
\item
Misalkan $(X,d)$ suatu ruang metrik dengan sekurang-kurangnya dua titik,
$d$ adalah metrik diskret, dan $p \in X$.
Hitunglah
$\operatorname{diam}(\{ p \})$ dan $\operatorname{diam}(X)$,
lalu simpulkan bahwa $(X,d)$ terbatas.
\end{enumerate}
\end{exercise}

\begin{exercise}
\pagebreak[2]
\leavevmode
\begin{enumerate}[a)]
\item
Carilah metrik $d$ pada $\N$ sedemikian sehingga $\N$ merupakan himpunan tak terbatas dalam $(\N,d)$.
\item
Carilah metrik $d$ pada $\N$ sedemikian sehingga $\N$ merupakan himpunan terbatas dalam $(\N,d)$.
\item
Carilah metrik $d$ pada $\N$ sedemikian sehingga untuk setiap $n \in \N$ dan setiap $\epsilon > 0$,
terdapat $m \in \N$, $m \neq n$, sedemikian sehingga $d(n,m) < \epsilon$.
\end{enumerate}
\end{exercise}

\begin{exercise} \label{exercise:C1ab}
Misalkan $C^1\bigl([a,b],\R\bigr)$\glsadd{not:contdifffuncs} himpunan fungsi yang memiliki turunan pertama
kontinu pada $[a,b]$.
Definisikan
\begin{equation*}
d(f,g) \coloneqq \snorm{f-g}_{[a,b]} + \snorm{f'-g'}_{[a,b]},
\end{equation*}
di mana $\snorm{\cdot}_{[a,b]}$ adalah norma seragam.  Buktikan bahwa $d$ merupakan metrik.
\end{exercise}

\begin{samepage}
\begin{exercise}
Perhatikan $\ell^2$, yakni himpunan barisan $\{ x_n \}_{n=1}^\infty$
bilangan real
sedemikian sehingga $\sum_{n=1}^\infty x_n^2 < \infty$.
\begin{enumerate}[a)]
\item
Buktikan \myindex{ketaksamaan Cauchy--Schwarz} untuk dua barisan
$\{x_n \}_{n=1}^\infty$ dan $\{ y_n \}_{n=1}^\infty$ dalam $\ell^2$:  Buktikan bahwa
$\sum_{n=1}^\infty x_n y_n$ konvergen (secara mutlak) dan
\begin{equation*}
{\biggl( \sum_{n=1}^\infty x_n y_n \biggr)}^2
\leq
\biggl(\sum_{n=1}^\infty x_n^2 \biggr)
\biggl(\sum_{n=1}^\infty y_n^2 \biggr) .
\end{equation*}
\item
Buktikan bahwa $\ell^2$ merupakan ruang metrik dengan metrik
$d(x,y) \coloneqq \sqrt{\sum_{n=1}^\infty {(x_n-y_n)}^2}$.
Petunjuk: Jangan lupa menunjukkan bahwa deret untuk $d(x,y)$ selalu konvergen
ke suatu bilangan berhingga.
\end{enumerate}
\end{exercise}
\end{samepage}

%%%%%%%%%%%%%%%%%%%%%%%%%%%%%%%%%%%%%%%%%%%%%%%%%%%%%%%%%%%%%%%%%%%%%%%%%%%%%%

\sectionnewpage
\section{Himpunan terbuka dan tertutup}
\label{sec:mettop}

%mbxINTROSUBSECTION

\sectionnotes{2 kuliah}

\subsection{Topologi}

Sebelum membahas konvergensi,
kita mendefinisikan \emph{\myindex{topologi}}.
Artinya,
kita mendefinisikan pengertian himpunan terbuka dan tertutup.
Sebelumnya,
kita terlebih dahulu mendefinisikan dua himpunan khusus yang terbuka dan tertutup.

\begin{defn}
Misalkan $(X,d)$ ruang metrik, $x \in X$, dan $\delta > 0$.  Definisikan
\emph{\myindex{bola terbuka}}, atau cukup \emph{\myindex{bola}}, berjari-jari $\delta$
yang berpusat di $x$ sebagai
\glsadd{not:openball}
\begin{equation*}
B(x,\delta) \coloneqq \bigl\{ y \in X : d(x,y) < \delta \bigr\} .
\end{equation*}
Definisikan \emph{\myindex{bola tertutup}} sebagai
\glsadd{not:closedball}
\begin{equation*}
C(x,\delta) \coloneqq \bigl\{ y \in X : d(x,y) \leq \delta \bigr\} .
\end{equation*}
\end{defn}

Ketika berurusan dengan ruang-ruang metrik yang berbeda, terkadang
penting untuk menegaskan ruang metrik tempat bola tersebut berada.  Untuk itu, kita
menulis $B_X(x,\delta) \coloneqq B(x,\delta)$ atau $C_X(x,\delta) \coloneqq C(x,\delta)$.

\begin{example}
Ambil ruang metrik $\R$ dengan metrik standar.  Untuk
$x \in \R$ dan $\delta > 0$,
\begin{equation*}
B(x,\delta) = (x-\delta,x+\delta) \qquad \text{dan} \qquad
C(x,\delta) = [x-\delta,x+\delta] .
\end{equation*}
\end{example}

\begin{example}
Berhati-hatilah ketika bekerja pada subruang.  Tinjau
ruang metrik $[0,1]$ sebagai subruang $\R$.  Maka, dalam $[0,1]$,
\begin{equation*}
B(0,\nicefrac{1}{2}) = B_{[0,1]}(0,\nicefrac{1}{2})
= \bigl\{ y \in [0,1] : \sabs{0-y} < \nicefrac{1}{2} \bigr\}
= [0,\nicefrac{1}{2}) .
\end{equation*}
Hasil ini berbeda dengan $B_{\R}(0,\nicefrac{1}{2}) =
(\nicefrac{-1}{2},\nicefrac{1}{2})$.
Hal penting yang harus selalu diperhatikan adalah ruang metrik tempat kita
bekerja.
\end{example}

\begin{defn}
Misalkan $(X,d)$ ruang metrik.  Himpunan bagian $V \subset X$
disebut \emph{terbuka}\index{himpunan terbuka}
jika untuk setiap $x \in V$, terdapat $\delta > 0$ sedemikian sehingga
$B(x,\delta) \subset V$.  Lihat \figureref{fig:msopenset}.  Himpunan bagian $E \subset X$ disebut
\emph{tertutup}\index{himpunan tertutup} jika komplemennya, $E^c = X \setminus E$, terbuka.
Jika ruang semesta $X$ tidak jelas dari konteks,
kita mengatakan \emph{$V$ terbuka dalam $X$} dan \emph{$E$ tertutup dalam $X$}.

Jika $x \in V$ dan $V$ terbuka, kita mengatakan
$V$ merupakan \emph{\myindex{lingkungan terbuka}} bagi $x$ (atau
terkadang cukup \emph{\myindex{lingkungan}}).
\end{defn}

\begin{myfigureht}
\myincludepdft{msopenset}{%
Suatu daerah pada bidang ditandai dengan warna abu-abu muda,
berbatas titik-titik, dan berlabel V.
Di dalam daerah tersebut terdapat lingkaran bertitik-titik
yang berpusat pada titik berlabel x, berjari-jari delta, dan berlabel
B dari x koma delta.  Lingkaran beserta bagian dalamnya,
yang diarsir abu-abu lebih gelap, seluruhnya berada di dalam
daerah abu-abu muda.}
\caption{Himpunan terbuka dalam ruang metrik.  Perhatikan bahwa $\delta$ bergantung pada $x$.\label{fig:msopenset}}
\end{myfigureht}

Secara intuitif, himpunan terbuka $V$ adalah himpunan yang tidak memuat
\myquote{batasnya}.
Di mana pun kita berada dalam $V$,
kita dapat \myquote{bergerak} sedikit dan
tetap berada dalam $V$.  Demikian pula, himpunan $E$ tertutup jika setiap titik di luar $E$
berjarak positif dari $E$.
Bola terbuka dan bola tertutup masing-masing merupakan contoh himpunan terbuka dan tertutup
(hal ini masih harus dibuktikan).
Namun, tidak setiap himpunan bersifat terbuka atau tertutup.  Pada umumnya, kebanyakan himpunan bagian tidak memiliki kedua sifat tersebut.

\begin{example}
Himpunan $(0,\infty) \subset \R$ terbuka:  Untuk sebarang $x \in (0,\infty)$,
ambil $\delta \coloneqq x$.  Maka $B(x,\delta) = (0,2x) \subset (0,\infty)$.

Himpunan $[0,\infty) \subset \R$ tertutup:  Untuk $x \in (-\infty,0) =
[0,\infty)^c$,
ambil $\delta \coloneqq -x$.  Maka $B(x,\delta) = (2x,0) \subset
(-\infty,0) = [0,\infty)^c$.

Himpunan $[0,1) \subset \R$ tidak terbuka maupun tertutup.  Pertama,
setiap bola dalam $\R$ yang berpusat di $0$, yaitu $B(0,\delta) = (-\delta,\delta)$, memuat bilangan
negatif sehingga tidak termuat dalam $[0,1)$.  Jadi, $[0,1)$ tidak terbuka.
Kedua, setiap bola dalam $\R$ yang berpusat di $1$,
$B(1,\delta) = (1-\delta,1+\delta)$, memuat
bilangan yang lebih kecil dari 1 dan lebih besar dari 0
(misalnya,\ $1-\nicefrac{\delta}{2}$ selama $\delta < 2$).
Dengan demikian, $[0,1)^c = \R \setminus
[0,1)$ tidak terbuka, sehingga $[0,1)$ tidak tertutup.

Jika $(X,d)$ sebarang ruang metrik dan $x \in X$ suatu titik, maka $\{ x \}$
tertutup (latihan).
Di sisi lain, $\{ x \}$ mungkin terbuka atau tidak, bergantung pada $X$.
Himpunan $\{ 0 \} \subset \R$ tidak terbuka karena $B(0,\delta)$
memuat bilangan taknol untuk setiap $\delta > 0$.
Jika $X=\{ x \}$, maka $\{ x \}$ terbuka.
\end{example}

\begin{prop} \label{prop:topology:open}
Misalkan $(X,d)$ ruang metrik.
\begin{enumerate}[(i)]
\item \label{topology:openi} $\emptyset$ dan $X$ terbuka.
\item \label{topology:openii} Jika $V_1, V_2, \ldots, V_k$ merupakan himpunan bagian terbuka dari $X$, maka
\begin{equation*}
\bigcap_{j=1}^k V_j
\end{equation*}
juga terbuka.  Artinya, irisan berhingga himpunan-himpunan terbuka bersifat terbuka.
\item \label{topology:openiii} Jika $\{ V_\lambda \}_{\lambda \in I}$ merupakan
sebarang koleksi himpunan bagian terbuka dari $X$, maka
\begin{equation*}
\bigcup_{\lambda \in I} V_\lambda
\end{equation*}
juga terbuka.  Artinya, gabungan himpunan-himpunan terbuka bersifat terbuka.
\end{enumerate}
\end{prop}

Himpunan indeks $I$ dalam \ref{topology:openiii} boleh berukuran sebarang.
Dengan $\bigcup_{\lambda \in I} V_\lambda$, yang dimaksud hanyalah himpunan
semua~$x$ sedemikian sehingga $x \in V_\lambda$ untuk setidaknya satu $\lambda \in I$.

\begin{proof}
Himpunan $\emptyset$ dan $X$ jelas terbuka dalam $X$.

Mari buktikan \ref{topology:openii}.
Jika $x \in \bigcap_{j=1}^k V_j$, maka $x \in V_j$ untuk setiap $j$.
Karena semua $V_j$ terbuka, untuk setiap $j$ terdapat $\delta_j > 0$
sedemikian sehingga $B(x,\delta_j) \subset V_j$.  Ambil $\delta \coloneqq \min \{
\delta_1,\delta_2,\ldots,\delta_k \}$ dan perhatikan bahwa $\delta > 0$.  Kita memperoleh
$B(x,\delta) \subset B(x,\delta_j) \subset V_j$ untuk setiap $j$, sehingga
$B(x,\delta) \subset \bigcap_{j=1}^k V_j$.  Jadi, irisan tersebut terbuka.

Mari buktikan \ref{topology:openiii}.
Jika $x \in \bigcup_{\lambda \in I} V_\lambda$, maka $x \in V_\lambda$ untuk suatu
$\lambda \in I$.
Karena $V_\lambda$ terbuka, terdapat $\delta > 0$
sedemikian sehingga $B(x,\delta) \subset V_\lambda$.  Oleh karena itu,
$B(x,\delta) \subset \bigcup_{\lambda \in I} V_\lambda$,
sehingga gabungan tersebut terbuka.
\end{proof}

\begin{example}
Perhatikan perbedaan antara
butir
\ref{topology:openii} dan \ref{topology:openiii}.
Butir \ref{topology:openii} tidak berlaku bagi sebarang irisan.
Sebagai contoh,
$\bigcap_{n=1}^\infty (\nicefrac{-1}{n},\nicefrac{1}{n}) = \{ 0 \}$,
yang tidak terbuka.
\end{example}


Pembuktian proposisi berikut yang serupa untuk himpunan tertutup
diserahkan sebagai latihan.

\begin{prop} \label{prop:topology:closed}
%FIXMEevillayouthack
\pagebreak[2]
Misalkan $(X,d)$ ruang metrik.
\begin{enumerate}[(i)]
\item \label{topology:closedi} $\emptyset$ dan $X$ tertutup.
\item \label{topology:closedii} Jika $\{ E_\lambda \}_{\lambda \in I}$ merupakan
sebarang koleksi himpunan bagian tertutup dari $X$, maka
\begin{equation*}
\bigcap_{\lambda \in I} E_\lambda
\end{equation*}
juga tertutup.  Artinya, irisan himpunan-himpunan tertutup bersifat tertutup.
\item \label{topology:closediii} Jika $E_1, E_2, \ldots, E_k$ merupakan himpunan bagian
tertutup dari $X$, maka
\begin{equation*}
\bigcup_{j=1}^k E_j
\end{equation*}
juga tertutup.  Artinya, gabungan berhingga himpunan-himpunan tertutup bersifat tertutup.
\end{enumerate}
\end{prop}

Walaupun namanya demikian,
kita belum menunjukkan bahwa bola terbuka bersifat terbuka dan bola tertutup bersifat
tertutup.  Mari buktikan fakta-fakta ini untuk membenarkan istilah tersebut.

\begin{prop} \label{prop:topology:ballsopenclosed}
Misalkan $(X,d)$ ruang metrik, $x \in X$, dan $\delta > 0$.  Maka
$B(x,\delta)$ terbuka dan
$C(x,\delta)$ tertutup.
\end{prop}

\begin{proof}
Misalkan $y \in B(x,\delta)$.  Tetapkan $\alpha \coloneqq \delta-d(x,y)$.  Karena $\alpha
> 0$, tinjau $z \in B(y,\alpha)$.  Maka
\begin{equation*}
d(x,z) \leq d(x,y) + d(y,z) < d(x,y) + \alpha = d(x,y) + \delta-d(x,y) =
\delta .
\end{equation*}
Oleh karena itu, $z \in B(x,\delta)$ untuk setiap $z \in B(y,\alpha)$.  Jadi,
$B(y,\alpha) \subset B(x,\delta)$, sehingga
$B(x,\delta)$ terbuka.  Lihat \figureref{fig:ballisopen}.

\begin{myfigureht}
\myincludepdft{ballisopen}{%
Sebuah cakram berarsir dengan batas titik-titik, berjari-jari delta,
berpusat di x dan berlabel B dari x koma delta.
Cakram lain seluruhnya berada di dalam cakram yang lebih besar,
menyentuh batasnya dari dalam, berpusat
di y, dan berjari-jari alfa.
Titik ketiga, z, berada di dalam
cakram yang lebih kecil.  Titik x, y, dan z membentuk
segitiga bergaris putus-putus.}
\caption{Bukti bahwa $B(x,\delta)$ terbuka: $B(y,\alpha) \subset
B(x,\delta)$ dengan ilustrasi ketaksamaan segitiga.\label{fig:ballisopen}}
\end{myfigureht}

Pembuktian bahwa $C(x,\delta)$ tertutup diserahkan sebagai latihan.
\end{proof}

Sekali lagi, perhatikan ruang metrik tempat kita bekerja.
Himpunan $[0,\nicefrac{1}{2})$ merupakan
bola terbuka dalam $[0,1]$, sehingga $[0,\nicefrac{1}{2})$
terbuka dalam $[0,1]$.  Di sisi lain, $[0,\nicefrac{1}{2})$
tidak terbuka maupun tertutup dalam $\R$.

\begin{prop} \label{prop:topology:intervals:openclosed}
Misalkan $a, b$ dua bilangan real dengan $a < b$.  Maka $(a,b)$, $(a,\infty)$,
dan $(-\infty,b)$ terbuka dalam $\R$.
Selain itu, $[a,b]$, $[a,\infty)$,
dan $(-\infty,b]$ tertutup dalam $\R$.
\end{prop}

Pembuktiannya diserahkan sebagai latihan.  Ingat bahwa
masih banyak himpunan terbuka dan
tertutup lain dalam himpunan bilangan real. %$\R$.

\begin{prop} \label{prop:topology:subspaceopen}
Misalkan $(X,d)$ ruang metrik dan $Y \subset X$.  Maka $U \subset Y$
terbuka dalam $Y$ (dalam topologi subruang) jika dan hanya jika
terdapat himpunan terbuka $V \subset X$ (yakni terbuka dalam $X$) sedemikian sehingga
$V \cap Y = U$.
\end{prop}

Sebagai contoh, tetapkan $X \coloneqq \R$, $Y \coloneqq [0,1]$, dan $U \coloneqq [0,\nicefrac{1}{2})$.
Kita telah melihat bahwa $U$ terbuka dalam $Y$.
Kita dapat mengambil $V \coloneqq (\nicefrac{-1}{2},\nicefrac{1}{2})$.

\begin{proof}
Andaikan $V \subset X$ terbuka dan $V \cap Y = U$.
Misalkan $x \in U$.
Karena $V$ terbuka dan $x \in V$, terdapat
$\delta > 0$ sedemikian sehingga $B_X(x,\delta) \subset V$.
Maka
\begin{equation*}
B_Y(x,\delta) = B_X(x,\delta) \cap Y \subset V \cap Y = U .
\end{equation*}
Jadi, $U$ terbuka dalam $Y$.

Pembuktian arah sebaliknya, yaitu bahwa jika $U \subset Y$
terbuka dalam topologi subruang maka terdapat $V$ demikian, diserahkan sebagai
\exerciseref{exercise:mssubspace}.
\end{proof}

Petunjuk untuk menyelesaikan pembuktian tersebut (latihan) adalah bahwa
himpunan terbuka dapat dipandang sebagai gabungan bola-bola terbuka.  Jika $U$
terbuka, maka untuk setiap $x \in U$, terdapat $\delta_x > 0$ (bergantung pada $x$) sedemikian sehingga
$B(x,\delta_x) \subset U$.  Maka $U = \bigcup_{x\in U} B(x,\delta_x)$.

Untuk himpunan bagian terbuka dari himpunan terbuka atau himpunan bagian tertutup dari himpunan
tertutup, keadaannya lebih sederhana.

\begin{prop} \label{prop:topology:subspacesame}
Misalkan $(X,d)$ ruang metrik, $V \subset X$ terbuka,
dan $E \subset X$ tertutup.
\begin{enumerate}[(i)]
\item \label{prop:topology:subspacesame:i}
$U \subset V$ terbuka dalam topologi subruang jika dan hanya jika $U$ terbuka
dalam $X$.
\item \label{prop:topology:subspacesame:ii}
$F \subset E$ tertutup dalam topologi subruang jika dan hanya jika $F$
tertutup dalam $X$.
\end{enumerate}
\end{prop}

\begin{proof}
Kita buktikan
\ref{prop:topology:subspacesame:i}
dan menyerahkan
\ref{prop:topology:subspacesame:ii} sebagai latihan.

Jika $U \subset V$ terbuka dalam topologi subruang, menurut
\propref{prop:topology:subspaceopen}, terdapat himpunan $W \subset X$
yang terbuka dalam $X$ sedemikian sehingga $U = W \cap V$.  Irisan dua himpunan terbuka
bersifat terbuka, sehingga $U$ terbuka dalam $X$.

Sekarang andaikan $U$ terbuka dalam $X$.  Maka $U = U \cap V$.  Jadi,
$U$ terbuka dalam $V$, sekali lagi menurut \propref{prop:topology:subspaceopen}.
\end{proof}

\subsection{Himpunan terhubung}

Mari memperumum gagasan interval ke ruang metrik umum.  Salah satu
ciri utama interval dalam $\R$ adalah sifat
terhubungnya---kita dapat bergerak secara kontinu dari satu titik di dalamnya ke
titik lain tanpa melompat.
Sebagai contoh, dalam $\R$ kita biasanya mempelajari fungsi pada interval,
sedangkan dalam ruang metrik yang lebih umum kita biasanya mempelajari fungsi pada himpunan terhubung.

\begin{defn}
Ruang metrik $(X,d)$ yang tidak kosong%
\footnote{Sebagian penulis tidak mengecualikan himpunan kosong dari definisi,
sehingga himpunan kosong terhubung menurut definisi mereka.
Kita menghindari himpunan kosong pada dasarnya karena alasan yang sama dengan
1 bukan bilangan prima maupun komposit:  Himpunan terhubung kita memiliki tepat
dua himpunan bagian terbuka-tertutup, sedangkan himpunan tak terhubung memiliki lebih dari dua.
Himpunan kosong hanya memiliki satu.}
disebut \emph{\myindex{terhubung}} jika himpunan bagian dari $X$ yang sekaligus
terbuka dan tertutup (terkadang disebut himpunan bagian
\emph{\myindex{terbuka-tertutup}}) hanyalah $\emptyset$ dan $X$ itu sendiri.
Jika $(X,d)$ yang tidak kosong tidak terhubung, kita menyebutnya
\emph{\myindex{tak terhubung}}.

Ketika istilah \emph{terhubung} diterapkan pada himpunan bagian tidak kosong $A \subset X$,
yang dimaksud ialah bahwa $A$, dengan topologi subruangnya, terhubung.
\end{defn}

Dengan kata lain, $X$ yang tidak kosong terhubung jika setiap kali kita menulis
$X = X_1 \cup X_2$, dengan $X_1 \cap X_2 = \emptyset$ serta $X_1$ dan $X_2$
terbuka, maka $X_1 = \emptyset$ atau $X_2 = \emptyset$.
Jadi, untuk menunjukkan bahwa $X$ tak terhubung, kita harus menemukan himpunan terbuka
tidak kosong yang saling lepas, $X_1$ dan
$X_2$, yang gabungannya adalah $X$.
Untuk himpunan bagian, gagasan ini kita nyatakan sebagai proposisi.
Proposisi tersebut diilustrasikan dalam \figureref{fig:disconsubset}.

\begin{prop}
Misalkan $(X,d)$ ruang metrik.
Himpunan tidak kosong $S \subset X$ tak terhubung jika dan hanya jika
terdapat himpunan terbuka $U_1$ dan
$U_2$ dalam $X$ sedemikian sehingga $U_1 \cap U_2 \cap S = \emptyset$,
$U_1 \cap S \neq \emptyset$,
$U_2 \cap S \neq \emptyset$, dan
\begin{equation*}
S = 
\bigl( U_1 \cap S \bigr)
\cup
\bigl( U_2 \cap S \bigr) .
\end{equation*}
\end{prop}

\begin{myfigureht}
\myincludepdft{disconsubset}{%
Dua daerah berarsir dengan garis batas penuh sama-sama
berlabel S.  Sebuah garis titik-titik mengelilingi setiap daerah
berarsir.
Bagian dalam daerah pertama ditandai U subskrip 1,
dan bagian dalam daerah kedua ditandai U subskrip 2.
Bagian dalam kedua garis titik-titik tersebut saling beririsan, tetapi
setiap daerah berarsir hanya berada di dalam salah satu garis
titik-titik.}
\caption{Himpunan bagian tak terhubung.  Perhatikan bahwa $U_1 \cap U_2$ tidak harus
kosong, tetapi $U_1 \cap U_2 \cap S = \emptyset$.\label{fig:disconsubset}}
\end{myfigureht}


\begin{proof}
Pertama, andaikan $S$ tak terhubung: Terdapat
$S_1$ dan $S_2$ yang tidak kosong, saling lepas, dan
terbuka dalam $S$, serta $S = S_1 \cup S_2$.
\propref{prop:topology:subspaceopen} menyatakan bahwa terdapat $U_1$ dan $U_2$
yang terbuka dalam $X$ sedemikian sehingga $U_1 \cap S = S_1$ dan
$U_2 \cap S = S_2$.

Untuk arah sebaliknya, mulailah dengan $U_1$ dan $U_2$ tersebut.
Maka $U_1 \cap S$ dan $U_2 \cap S$ terbuka dalam $S$ menurut
\propref{prop:topology:subspaceopen}.
Menurut pembahasan sebelum proposisi, $S$ tak terhubung.
\end{proof}

\begin{example}
Misalkan $S \subset \R$ dan terdapat $x,y,z$ sedemikian sehingga $x < z < y$, dengan $x,y \in S$
dan $z \notin S$.  Klaim: \emph{$S$ tak terhubung}.  Bukti:  Perhatikan bahwa
\begin{equation*}
\bigl( (-\infty,z) \cap S \bigr)
\cup
\bigl( (z,\infty) \cap S \bigr)
= S .
\end{equation*}
\end{example}

\begin{prop}
Himpunan tidak kosong $S \subset \R$ terhubung jika dan hanya jika $S$
merupakan interval atau satu titik.
\end{prop}

\begin{proof}
Andaikan $S$ terhubung.  Jika $S$ hanya terdiri atas satu titik,
pembuktian selesai.  Jadi, andaikan $x < y$ dan $x,y \in S$.  Jika $z \in \R$
sedemikian sehingga $x < z < y$, maka $(-\infty,z) \cap S$ tidak kosong dan $(z,\infty) \cap
S$ tidak kosong.  Kedua himpunan itu saling lepas.  Karena
$S$ terhubung, gabungan keduanya tidak mungkin sama dengan $S$, sehingga $z \in S$.
Menurut \propref{prop:intervaldef}, $S$ merupakan interval.

Jika $S$ hanya terdiri atas satu titik, maka himpunan tersebut terhubung.
Selanjutnya, andaikan $S$ merupakan interval.
Tinjau himpunan bagian terbuka $U_1$ dan $U_2$ dari $\R$ sedemikian sehingga
$U_1 \cap S$ dan $U_2 \cap S$ tidak kosong, dan
$S = 
\bigl( U_1 \cap S \bigr)
\cup
\bigl( U_2 \cap S \bigr)$.  Kita akan menunjukkan bahwa $U_1 \cap S$
dan $U_2 \cap S$ memiliki suatu titik bersama, sehingga keduanya tidak saling lepas;
hal ini membuktikan bahwa $S$ terhubung.
Misalkan $x \in U_1 \cap S$
dan $y \in U_2 \cap S$.  Tanpa mengurangi keumuman, andaikan $x < y$.
Karena $S$ merupakan interval, $[x,y] \subset S$.
Perhatikan bahwa $U_2 \cap [x,y] \neq \emptyset$, dan
tetapkan $z \coloneqq \inf (U_2 \cap [x,y])$.
Kita ingin menunjukkan bahwa $z \in U_1$.
Jika $z = x$, maka $z \in U_1$.
Jika $z > x$,
maka untuk setiap $\epsilon > 0$, bola $B(z,\epsilon) =
(z-\epsilon,z+\epsilon)$ memuat titik-titik dalam $[x,y]$ yang tidak berada dalam $U_2$,
karena $z$ merupakan infimum $U_2 \cap [x,y]$.
Jadi, $z \notin U_2$ karena $U_2$ terbuka.
Oleh karena itu, $z \in U_1$ sebab setiap titik dalam $[x,y]$ berada dalam
$U_1$ atau $U_2$.
Karena $U_1$ terbuka,
$B(z,\delta) \subset U_1$ untuk $\delta > 0$ yang cukup kecil.
Karena $z$ merupakan infimum himpunan tidak kosong $U_2 \cap [x,y]$,
harus terdapat suatu $w \in U_2 \cap [x,y]$
sedemikian sehingga $w \in [z,z+\delta) \subset B(z,\delta) \subset U_1$.
Oleh karena itu, $w \in U_1 \cap U_2 \cap [x,y]$.
Jadi, $U_1 \cap S$ dan $U_2 \cap S$ tidak saling lepas, sehingga
$S$ terhubung.
Lihat \figureref{fig:intervalcon}.
\end{proof}

\begin{myfigureht}
\myincludepdft{intervalcon}{%
Diagram garis bilangan dengan titik
x yang ditandai di sebelah kiri dan titik y di sebelah
kanan.  Di antara x dan y terdapat
titik z dan w.  Suatu interval terbuka
yang bermula di suatu tempat di sebelah kiri dan berlanjut melewati w (sehingga
memuat x, z, dan w) ditandai U subskrip 1.
Suatu interval terbuka dari z yang berlanjut keluar dari diagram
ke kanan ditandai U subskrip 2.
Sebuah interval terbuka kecil dari z dikurangi delta
hingga z ditambah delta yang memuat w juga ditandai.}
\caption{Bukti bahwa suatu interval terhubung.\label{fig:intervalcon}}
\end{myfigureht}

\begin{example}
Sering kali sebuah bola $B(x,\delta)$ terhubung, tetapi hal ini tidak
selalu benar dalam setiap ruang metrik.
Sebagai contoh paling sederhana,
ambil ruang dua titik $\{ a, b\}$ dengan metrik diskret.
Maka $B(a,2) = \{ a , b \}$ tidak
terhubung karena $B(a,1) = \{ a \}$ dan
$B(b,1) = \{ b \}$ terbuka dan saling lepas.
\end{example}

\subsection{Penutupan dan batas}

Terkadang kita ingin mengambil suatu himpunan dan menambahkan semua titik yang dapat didekati
dari dalam himpunan tersebut.  Konsep ini disebut penutupan.

\begin{defn}
Misalkan $(X,d)$ ruang metrik dan $A \subset X$.
Kita mendefinisikan \emph{\myindex{penutupan}} dari $A$ sebagai himpunan
\glsadd{not:closure}
\begin{equation*}
\widebar{A} \coloneqq \bigcap \{ E \subset X : E \text{ tertutup dan } A \subset
E \} .
\end{equation*}
\end{defn}

Artinya,
$\widebar{A}$ merupakan irisan semua himpunan tertutup yang memuat $A$.


\begin{prop}
Misalkan $(X,d)$ ruang metrik dan $A \subset X$.  Penutupan $\widebar{A}$
tertutup dan $A \subset \widebar{A}$.
Selain itu, jika $A$ tertutup, maka $\widebar{A} = A$.
\end{prop}

\begin{proof}
Penutupan merupakan irisan himpunan-himpunan tertutup, sehingga $\widebar{A}$ tertutup.
Terdapat setidaknya satu himpunan tertutup yang memuat $A$, yaitu $X$ sendiri,
sehingga $A \subset \widebar{A}$.
Jika $A$ tertutup, maka $A$ adalah himpunan tertutup yang memuat $A$.
Jadi, $\widebar{A} \subset A$, sehingga $A = \widebar{A}$.
\end{proof}

\begin{example}
Penutupan $(0,1)$ dalam $\R$ adalah $[0,1]$.  Bukti:  Jika
$E$ tertutup dan memuat $(0,1)$, maka $E$ memuat $0$ dan $1$ (mengapa?).
Dengan demikian, $[0,1] \subset E$.  Namun, $[0,1]$ juga tertutup.
Jadi, penutupannya adalah $\overline{(0,1)} = [0,1]$.
\end{example}

\begin{example}
Perhatikan dengan cermat ruang metrik tempat kita bekerja.
Jika $X = (0,\infty)$, maka
penutupan $(0,1)$ dalam $(0,\infty)$ adalah $(0,1]$.  Bukti:  Seperti
di atas, $(0,1]$ tertutup dalam $(0,\infty)$ (mengapa?).  Setiap himpunan tertutup $E$
yang memuat $(0,1)$ harus memuat 1 (mengapa?).  Oleh karena itu, $(0,1] \subset E$,
sehingga $\overline{(0,1)} = (0,1]$ ketika kita bekerja dalam $(0,\infty)$.
\end{example}

Mari kita buktikan bahwa penutupan memuat semua titik yang dapat kita
\myquote{dekati} dari dalam himpunan.

\begin{prop} \label{prop:msclosureappr}
Misalkan $(X,d)$ ruang metrik dan $A \subset X$.  Maka $x \in \widebar{A}$
jika dan hanya jika untuk setiap $\delta > 0$, berlaku $B(x,\delta) \cap A \neq \emptyset$.
\end{prop}

\begin{proof}
Mari buktikan kedua kontraposisinya.
Kita akan menunjukkan bahwa $x \notin \widebar{A}$ jika dan hanya jika terdapat
$\delta > 0$ sedemikian sehingga $B(x,\delta) \cap A = \emptyset$.

Pertama, andaikan $x \notin \widebar{A}$.  Kita mengetahui bahwa $\widebar{A}$
tertutup.  Dengan demikian, terdapat $\delta > 0$ sedemikian sehingga
$B(x,\delta) \subset \widebar{A}^c$.  Karena $A \subset \widebar{A}$, kita
memperoleh $B(x,\delta) \subset \widebar{A}^c \subset A^c$, sehingga
$B(x,\delta) \cap A = \emptyset$.

Sebaliknya, andaikan terdapat $\delta > 0$ sedemikian sehingga
$B(x,\delta) \cap A = \emptyset$.
Dengan kata lain,
$A \subset {B(x,\delta)}^c$.
Karena
${B(x,\delta)}^c$ merupakan himpunan tertutup, $x \not \in {B(x,\delta)}^c$,
dan $\widebar{A}$ merupakan irisan
himpunan-himpunan tertutup yang memuat $A$, kita memperoleh $x \notin \widebar{A}$.
\end{proof}

Kita juga dapat membahas bagian dalam suatu himpunan
(titik-titik yang tidak dapat kita dekati dari komplemennya),
serta batas suatu himpunan (titik-titik yang dapat kita
dekati baik dari himpunan maupun dari komplemennya).

\begin{defn}
Misalkan $(X,d)$ ruang metrik dan $A \subset X$.
Kita mendefinisikan \emph{\myindex{bagian dalam}} dari $A$ sebagai himpunan
\glsadd{not:interior}
\begin{equation*}
A^\circ \coloneqq \{ x \in A : \text{terdapat } \delta > 0
\text{ sedemikian sehingga } B(x,\delta) \subset A \} .
\end{equation*}
Kita mendefinisikan \emph{\myindex{batas}} dari $A$ sebagai himpunan
\glsadd{not:boundary}
\begin{equation*}
\partial A \coloneqq \widebar{A}\setminus A^\circ.
\end{equation*}
\end{defn}

Sebagai rumusan lain, bagian dalam merupakan gabungan himpunan-himpunan terbuka yang termuat dalam $A$;
lihat \exerciseref{exercise:interiortopologic}.  Berdasarkan
definisi, $A^\circ \subset A$;
namun, titik-titik batas dapat berada atau tidak berada dalam $A$.

\begin{example}
Misalkan $A \coloneqq (0,1]$ dan $X \coloneqq \R$.  Maka $\widebar{A}=[0,1]$, $A^\circ = (0,1)$,
dan $\partial A = \{ 0, 1 \}$.
\end{example}

\begin{example}
Tinjau $X \coloneqq \{ a, b \}$ dengan metrik diskret,
dan tetapkan $A \coloneqq \{ a \}$.  Maka $\widebar{A} = A^\circ = A$ dan $\partial A =
\emptyset$.
\end{example}


\begin{prop}
Misalkan $(X,d)$ ruang metrik dan $A \subset X$.  Maka $A^\circ$ terbuka
dan $\partial A$ tertutup.
\end{prop}

\begin{proof}
Untuk $x \in A^\circ$, terdapat $\delta > 0$ sedemikian sehingga $B(x,\delta)
\subset A$.  Jika $z \in B(x,\delta)$, karena bola terbuka bersifat terbuka,
terdapat $\epsilon > 0$ sedemikian sehingga $B(z,\epsilon) \subset B(x,\delta)
\subset A$.  Jadi, $z \in A^\circ$.  Oleh karena itu, $B(x,\delta) \subset
A^\circ$, sehingga $A^\circ$ terbuka.

Karena $A^\circ$ terbuka,
$\partial A = \widebar{A} \setminus A^\circ = \widebar{A} \cap
{(A^\circ)}^c$ tertutup.
\end{proof}

Batas adalah himpunan titik-titik yang dekat dengan himpunan maupun
komplemennya.  Lihat \figureref{fig:msboundary} untuk diagram
proposisi berikut.

\begin{prop}
Misalkan $(X,d)$ ruang metrik dan $A \subset X$.  Maka $x \in \partial A$
jika dan hanya jika untuk setiap $\delta > 0$,
$B(x,\delta) \cap A$ dan
$B(x,\delta) \cap A^c$ keduanya tidak kosong.
\end{prop}

\begin{myfigureht}
\myincludepdft{msboundary}{%
Gambar di sekitar batas himpunan A yang diarsir.
Batas digambarkan sebagai garis penuh, dengan A di bawah
dan komplemen A di atasnya.  Lingkaran bertitik-titik
berjari-jari delta berpusat pada titik x di
batas dan ditandai sebagai B dari x koma delta.}
\caption{Batas adalah himpunan tempat setiap bola memuat titik dari himpunan
dan juga dari komplemennya.\label{fig:msboundary}}
\end{myfigureht}

\begin{proof}
Andaikan $x \in \partial A =  \widebar{A} \setminus A^\circ$ dan
ambil sebarang $\delta > 0$.
Menurut \propref{prop:msclosureappr}, $B(x,\delta)$ memuat
suatu titik dari $A$.  Jika $B(x,\delta)$ tidak memuat titik dari $A^c$,
maka $x$ berada dalam $A^\circ$.  Oleh karena itu, $B(x,\delta)$ juga memuat suatu titik dari
$A^c$.

Mari buktikan arah sebaliknya dengan kontraposisi.
Andaikan $x \notin \partial A$, sehingga $x \notin \widebar{A}$ atau $x \in A^\circ$.
Jika $x \notin \widebar{A}$, maka
$B(x,\delta) \subset \widebar{A}^c$
untuk suatu $\delta > 0$ karena $\widebar{A}$ tertutup.
Jadi, $B(x,\delta) \cap A$ kosong karena $\widebar{A}^c \subset
A^c$.
Jika $x \in A^\circ$, maka
$B(x,\delta) \subset A$ untuk suatu $\delta > 0$,
sehingga $B(x,\delta) \cap A^c$ kosong.
\end{proof}

Kita memperoleh akibat langsung berikut tentang penutupan $A$ dan $A^c$.  Kita
cukup menerapkan \propref{prop:msclosureappr}.

\begin{cor}
Misalkan $(X,d)$ ruang metrik dan $A \subset X$.
Maka $\partial A = \widebar{A} \cap \widebar{A^c}$.
\end{cor}

\subsection{Latihan}

\begin{exercise}
Buktikan \propref{prop:topology:closed}.  Petunjuk:
Terapkan \propref{prop:topology:open} pada
komplemen himpunan-himpunan tersebut.
\end{exercise}

\begin{exercise}
Selesaikan pembuktian \propref{prop:topology:ballsopenclosed} dengan
membuktikan bahwa $C(x,\delta)$ tertutup.
\end{exercise}

\begin{exercise}
Buktikan \propref{prop:topology:intervals:openclosed}.
\end{exercise}

\begin{exercise}
Misalkan $(X,d)$ ruang metrik tak kosong dengan topologi diskret.  Tunjukkan
bahwa $X$ terhubung jika dan hanya jika memuat tepat satu elemen.
\end{exercise}

\begin{exercise}
Ambil $\Q$ dengan metrik standar, $d(x,y) = \sabs{x-y}$, sebagai ruang metrik kita.
Buktikan bahwa $\Q$
\emph{\myindex{terpisah total}}, yaitu, tunjukkan 
bahwa untuk setiap $x, y \in \Q$ dengan $x \neq y$, terdapat
dua himpunan terbuka $U$ dan $V$ sedemikian sehingga $x \in U$, $y \in V$,
$U \cap V = \emptyset$, dan $U \cup V = \Q$.
\end{exercise}

\begin{exercise}
Tunjukkan bahwa dalam ruang metrik,
setiap himpunan terbuka dapat ditulis sebagai gabungan himpunan-himpunan tertutup.
\end{exercise}

\begin{samepage}
\begin{exercise}
Buktikan bahwa dalam ruang metrik,
\begin{enumerate}[a)]
\item
$E$ tertutup jika dan hanya jika $\partial E \subset E$.
\item
$U$ terbuka jika dan hanya jika $\partial U \cap U = \emptyset$.
\end{enumerate}
\end{exercise}
\end{samepage}

\begin{exercise}
Buktikan bahwa dalam ruang metrik,
\begin{enumerate}[a)]
\item
$A$ terbuka jika dan hanya jika $A^\circ = A$.
\item
$U \subset A^\circ$
untuk setiap himpunan terbuka $U$ sedemikian sehingga $U \subset A$.
\end{enumerate}
\end{exercise}

\begin{exercise}
Misalkan $X$ suatu himpunan serta $d$ dan $d'$ dua metrik pada $X$.
Andaikan terdapat $\alpha > 0$ dan $\beta > 0$
sedemikian sehingga $\alpha d(x,y) \leq d'(x,y) \leq \beta d(x,y)$ untuk semua $x,y \in X$.
Tunjukkan bahwa $U$ terbuka dalam $(X,d)$ jika dan hanya jika $U$ terbuka dalam $(X,d')$.
Artinya, topologi $(X,d)$ dan $(X,d')$ sama.
\end{exercise}


\begin{exercise}
Misalkan $\{ S_i \}$, $i \in \N$,
suatu koleksi subhimpunan terhubung dari ruang metrik $(X,d)$,
dan terdapat $x \in X$ sedemikian sehingga $x \in S_i$ untuk semua $i \in \N$.
Tunjukkan bahwa $\bigcup_{i=1}^\infty S_i$ terhubung.
\end{exercise}

\begin{exercise}
\pagebreak[2]
Misalkan $A$ himpunan terhubung dalam suatu ruang metrik.
\begin{enumerate}[a)]
\item
Apakah $\widebar{A}$ terhubung?  Buktikan atau temukan contoh penyangkal.
\item
Apakah $A^\circ$ terhubung?  Buktikan atau temukan contoh penyangkal.
\end{enumerate}
Petunjuk: Pikirkan himpunan-himpunan dalam $\R^2$.
\end{exercise}

\begin{exercise} \label{exercise:mssubspace}
Selesaikan pembuktian \propref{prop:topology:subspaceopen}.
Misalkan $(X,d)$ ruang metrik dan $Y \subset X$.  Tunjukkan bahwa
dengan metrik subruang pada $Y$, jika himpunan $U \subset Y$
terbuka (dalam $Y$), maka terdapat himpunan terbuka $V \subset X$ sedemikian
sehingga $U = V \cap Y$.
\end{exercise}

\begin{samepage}
\begin{exercise}
Misalkan $(X,d)$ ruang metrik.
\begin{enumerate}[a)]
\item
Untuk setiap $x \in X$ dan $\delta > 0$, tunjukkan
$\overline{B(x,\delta)} \subset C(x,\delta)$.
\item
Apakah selalu benar bahwa
$\overline{B(x,\delta)} = C(x,\delta)$?  Buktikan atau temukan contoh penyangkal.
\end{enumerate}
\end{exercise}
\end{samepage}

\begin{exercise} \label{exercise:interiortopologic}
Misalkan $(X,d)$ ruang metrik dan $A \subset X$.  Tunjukkan bahwa
$A^\circ = \bigcup \{ V : V \text{ terbuka dan } V \subset A \}$.
\end{exercise}

\begin{exercise}
Selesaikan pembuktian \propref{prop:topology:subspacesame}.
\end{exercise}

\begin{exercise}
Misalkan $(X,d)$ ruang metrik.  Tunjukkan bahwa terdapat metrik terbatas
$d'$ sedemikian sehingga $(X,d')$ memiliki himpunan-himpunan terbuka yang sama; dengan kata lain,
topologinya sama.
\end{exercise}

\begin{samepage}
\begin{exercise}
Misalkan $(X,d)$ ruang metrik.
\begin{enumerate}[a)]
\item
Buktikan bahwa untuk setiap $x \in X$, terdapat $\delta > 0$ sedemikian sehingga
$B(x,\delta) = \{ x \}$, atau $B(x,\delta)$ tak hingga untuk setiap $\delta >  0$.
\item
Temukan contoh eksplisit $(X,d)$ dengan $X$ tak hingga, sedemikian sehingga
untuk setiap $\delta > 0$ dan
setiap $x \in X$,
bola $B(x,\delta)$ berhingga.
\item
Temukan contoh eksplisit $(X,d)$ sedemikian sehingga untuk setiap $\delta > 0$ dan
setiap $x \in X$,
bola $B(x,\delta)$ tak hingga terhitung.
\item
Buktikan bahwa jika $X$ tak terhitung, maka terdapat
$x \in X$ dan $\delta > 0$ sedemikian sehingga $B(x,\delta)$ tak terhitung.
\end{enumerate}
\end{exercise}
\end{samepage}

\begin{exercise}
Untuk setiap $x \in \R^n$ dan setiap $\delta > 0$, definisikan \myquote{persegi panjang}
$R(x,\delta) \coloneqq
(x_1-\delta,x_1+\delta) \times
(x_2-\delta,x_2+\delta) \times \cdots \times
(x_n-\delta,x_n+\delta)$.  Tunjukkan bahwa himpunan-himpunan ini dan bola-bola
dalam metrik standar membangkitkan koleksi himpunan terbuka yang sama.  Artinya, tunjukkan bahwa himpunan $U \subset \R^n$
terbuka dalam pengertian metrik standar jika dan hanya jika untuk setiap
titik $x \in U$, terdapat $\delta > 0$ sedemikian sehingga $R(x,\delta) \subset
U$.
\end{exercise}

%%%%%%%%%%%%%%%%%%%%%%%%%%%%%%%%%%%%%%%%%%%%%%%%%%%%%%%%%%%%%%%%%%%%%%%%%%%%%%

\sectionnewpage
\section{Barisan dan kekonvergenan}
\label{sec:metseqs}

%mbxINTROSUBSECTION

\sectionnotes{1 kuliah}

\subsection{Barisan}

Pengertian barisan dalam ruang metrik sangat mirip dengan barisan
bilangan real.
Definisi-definisi terkait pada dasarnya sama dengan definisi untuk bilangan real
dalam \chapterref{seq:chapter}, dengan $\R$ yang dilengkapi
dengan metrik standar $d(x,y)=\sabs{x-y}$ digantikan
oleh ruang metrik sebarang $(X,d)$.


\begin{defn}
Suatu \emph{\myindex{barisan}} dalam ruang metrik $(X,d)$ adalah fungsi
$x \colon \N \to X$.  Seperti sebelumnya, kita menuliskan $x_n$ untuk elemen ke-$n$ dalam
barisan tersebut, dan menggunakan notasi berikut untuk seluruh barisan:
\glsadd{not:sequence}
\begin{equation*}
\{ x_n \}_{n=1}^\infty .
\end{equation*}

Barisan $\{ x_n \}_{n=1}^\infty$ disebut \emph{terbatas}\index{barisan terbatas} jika
terdapat titik $p \in X$ dan $B \in \R$ sedemikian sehingga
\begin{equation*}
d(p,x_n) \leq B \qquad \text{untuk semua } n \in \N.
\end{equation*}
Dengan kata lain, barisan $\{x_n\}_{n=1}^\infty$ terbatas apabila
himpunan $\{ x_n : n \in \N \}$
terbatas.

Jika $\{ n_k \}_{k=1}^\infty$ merupakan barisan bilangan asli
sedemikian sehingga $n_{k+1} > n_k$ untuk semua $k$, maka
barisan $\{ x_{n_k} \}_{k=1}^\infty$\glsadd{not:subsequence}
disebut
\emph{\myindex{subbarisan}} dari $\{x_n \}_{n=1}^\infty$.
\end{defn}

Dengan cara serupa, kita mendefinisikan kekonvergenan.
Lihat \figureref{fig:sequence-convergence-metric} untuk memperoleh gambaran tentang
definisi tersebut.

\begin{defn}
Barisan $\{ x_n \}_{n=1}^\infty$ dalam ruang metrik $(X,d)$ dikatakan
\emph{konvergen} ke titik
$p \in X$ jika untuk setiap $\epsilon > 0$, terdapat $M \in \N$ sedemikian
sehingga $d(x_n,p) < \epsilon$ untuk semua $n \geq M$.  Titik $p$
disebut \emph{limit}\index{limit!barisan dalam ruang metrik}
dari $\{ x_n \}_{n=1}^\infty$.
Kita akan membuktikan bahwa limit tersebut tunggal, sehingga
kita menuliskan\glsadd{not:limseq}
\begin{equation*}
\lim_{n\to \infty} x_n \coloneqq p .
\end{equation*}

Barisan
yang konvergen disebut \emph{konvergen}\index{konvergen!barisan dalam ruang
metrik}.
Jika tidak, barisan tersebut disebut 
\emph{divergen}\index{divergen!barisan dalam ruang metrik}.
\begin{myfigureht}
\myincludegraphics{sequence-convergence-metric}{%
Sebuah lingkaran bertitik-titik berjari-jari epsilon yang berpusat di p digambarkan.
Sejumlah titik dari awal barisan, dari x subskrip 1 hingga
x subskrip 10, ditampilkan.  Titik-titik x subskrip 1,
x subskrip 3, x subskrip 4, x subskrip 5, dan x subskrip 6 berada
di luar lingkaran, sedangkan semua titik lainnya berada di dalamnya.}
\caption{Barisan yang konvergen ke $p$.  Sepuluh titik pertama 
ditampilkan dan $M=7$ untuk $\epsilon$ ini.\label{fig:sequence-convergence-metric}}
\end{myfigureht}
\end{defn}

Kita mulai dengan membuktikan bahwa limit tersebut tunggal.  Pembuktiannya hampir
identik (kata demi kata) dengan pembuktian fakta yang sama untuk barisan
bilangan real,
\propref{prop:limisunique}.
Pembuktian banyak hasil yang telah kita ketahui tentang barisan bilangan real dapat
disesuaikan dengan
konteks ruang metrik yang lebih umum.  Kita harus mengganti $\sabs{x-y}$
dengan $d(x,y)$ dalam pembuktian dan menerapkan ketaksamaan segitiga dengan tepat.

\begin{prop} \label{prop:mslimisunique}
Barisan konvergen dalam ruang metrik memiliki limit tunggal.
\end{prop}

\begin{proof}
%NOTE: should be word for word the same as 2.1.6
Misalkan $\{ x_n \}_{n=1}^\infty$ memiliki limit $x$ dan $y$.
Ambil sebarang $\epsilon > 0$.
Dari definisi, pilih $M_1$ sedemikian sehingga untuk semua $n \geq M_1$,
$d(x_n,x) < \nicefrac{\epsilon}{2}$.  Dengan cara serupa, pilih $M_2$
sedemikian sehingga untuk semua $n \geq M_2$, berlaku
$d(x_n,y) < \nicefrac{\epsilon}{2}$.  Sekarang ambil $n$ sedemikian sehingga
$n \geq M_1$ dan juga $n \geq M_2$; maka diperoleh estimasi berikut:
\begin{equation*}
\begin{split}
d(y,x)
& \leq
d(y,x_n) + d(x_n,x) \\
& <
\frac{\epsilon}{2} + \frac{\epsilon}{2} = \epsilon .
\end{split}
\end{equation*}
Karena $d(y,x) < \epsilon$ untuk semua $\epsilon > 0$, kita memperoleh $d(x,y) = 0$
sehingga $y=x$.  Jadi, limit tersebut (jika ada) tunggal.
\end{proof}

Pembuktian proposisi-proposisi berikut diserahkan sebagai latihan.

\begin{prop} \label{prop:msconvbound}
Barisan konvergen dalam ruang metrik bersifat terbatas.
\end{prop}

\begin{prop} \label{prop:msconvifa}
Barisan $\{ x_n \}_{n=1}^\infty$ dalam ruang metrik $(X,d)$ konvergen ke $p \in X$
jika dan hanya
jika terdapat barisan bilangan real $\{ a_n \}_{n=1}^\infty$ sedemikian sehingga
\begin{equation*}
d(x_n,p) \leq a_n \quad \text{untuk semua } n \in \N,
\qquad
\text{dan}
\qquad
\lim_{n\to\infty} a_n = 0.
\end{equation*}
\end{prop}

\begin{prop} \label{prop:mssubseq}
Misalkan $\{ x_n \}_{n=1}^\infty$ barisan dalam ruang metrik $(X,d)$.
\begin{enumerate}[(i)]
\item Jika $\{ x_n \}_{n=1}^\infty$ konvergen ke $p \in X$, maka setiap
subbarisan $\{ x_{n_k} \}_{k=1}^\infty$
konvergen ke $p$.
\item Jika untuk suatu $K \in \N$, ekor-$K$ $\{ x_n \}_{n=K+1}^\infty$
konvergen ke $p \in X$, maka
 $\{ x_n \}_{n=1}^\infty$ konvergen ke $p$.
\end{enumerate}
\end{prop}

\begin{example}
Ambil $C\bigl([a,b],\R\bigr)$ sebagai himpunan fungsi kontinu dengan metrik
yang diinduksi oleh norma seragam.  Kita telah melihat bahwa ini menghasilkan ruang metrik.
Jika kita mencermati definisi kekonvergenan, kita melihat bahwa definisi tersebut identik
dengan konvergensi seragam.  Artinya, $\{ f_n \}_{n=1}^\infty$ konvergen seragam jika dan hanya
jika barisan tersebut konvergen dalam pengertian ruang metrik.
\end{example}

\begin{remark}
Mungkin mengejutkan bahwa pada himpunan fungsi $f \colon [a,b] \to
\R$ (kontinu ataupun tidak)
tidak ada metrik yang menghasilkan konvergensi titik demi titik karena domain $[a,b]$
tak terhitung.  Pembuktian
fakta ini berada di luar cakupan buku ini.
\end{remark}

\subsection{Konvergensi dalam ruang Euklides}

Dalam ruang Euklides $\R^n$, suatu barisan konvergen jika dan hanya jika setiap
komponennya konvergen:

\begin{prop} \label{prop:msconveuc}
Misalkan $\{ x_m \}_{m=1}^\infty$ suatu barisan dalam $\R^n$,
dengan $x_m = \bigl(x_{m,1},x_{m,2},\ldots,x_{m,n}\bigr) \in \R^n$.
Maka $\{ x_m \}_{m=1}^\infty$ konvergen jika dan hanya jika
$\{ x_{m,k} \}_{m=1}^\infty$ konvergen untuk setiap $k=1,2,\ldots,n$; dalam hal ini,
\begin{equation*}
\lim_{m\to\infty}
x_m =
\Bigl(
\lim_{m\to\infty} x_{m,1},
\lim_{m\to\infty} x_{m,2},
\ldots,
\lim_{m\to\infty} x_{m,n}
\Bigr) .
\end{equation*}
\end{prop}

\begin{proof}
Misalkan
$\{ x_m \}_{m=1}^\infty$ konvergen ke
$y = (y_1,y_2,\ldots,y_n) \in \R^n$.
Untuk $\epsilon > 0$, terdapat $M$ sedemikian sehingga untuk semua
$m \geq M$, berlaku
\begin{equation*}
d(y,x_m) < \epsilon.
\end{equation*}
Tetapkan sebarang $k=1,2,\ldots,n$.  Untuk semua $m \geq M$,
\begin{equation*}
\bigl\lvert y_k - x_{m,k} \bigr\rvert
=
\sqrt{{\bigl(y_k - x_{m,k} \bigr)}^2}
\leq
\sqrt{\sum_{\ell=1}^n {\bigl(y_\ell-x_{m,\ell}\bigr)}^2}
= d(y,x_m) < \epsilon .
\end{equation*}
Jadi, barisan $\{ x_{m,k} \}_{m=1}^\infty$ konvergen ke $y_k$.

Untuk arah sebaliknya, misalkan 
$\{ x_{m,k} \}_{m=1}^\infty$ konvergen ke $y_k$ untuk setiap $k=1,2,\ldots,n$.
Untuk $\epsilon > 0$, pilih $M$ sedemikian sehingga jika $m \geq M$, maka 
$\bigl\lvert y_k-x_{m,k} \bigr\rvert < \nicefrac{\epsilon}{\sqrt{n}}$ untuk semua
$k=1,2,\ldots,n$.  Maka
\begin{equation*}
d(y,x_m)
=
\sqrt{\sum_{k=1}^n {\bigl(y_k-x_{m,k}\bigr)}^2}
<
\sqrt{\sum_{k=1}^n {\left(\frac{\epsilon}{\sqrt{n}}\right)}^2}
=
\sqrt{\sum_{k=1}^n \frac{{\epsilon^2}}{n}}
= \epsilon .
\end{equation*}
Artinya, barisan $\{ x_m \}_{m=1}^\infty$ konvergen ke
$y = (y_1,y_2,\ldots,y_n) \in \R^n$.
\end{proof}

\begin{example}
Seperti telah disebutkan, himpunan bilangan kompleks $\C$, dengan $z = x+iy$, dipandang 
sebagai ruang metrik $\R^2$.  Proposisi tersebut menyatakan bahwa
barisan $\{ z_n \}_{n=1}^\infty = \{ x_n + iy_n \}_{n=1}^\infty$ konvergen
ke $z = x+iy$
jika dan hanya jika $\{ x_n \}_{n=1}^\infty$ konvergen ke $x$ dan 
$\{ y_n \}_{n=1}^\infty$ konvergen ke $y$.
\end{example}

\subsection{Kekonvergenan dan topologi}

Topologi---himpunan semua himpunan terbuka dari suatu ruang---menentukan
barisan mana yang konvergen.

\begin{prop} \label{prop:msconvtopo}
Misalkan $(X,d)$ ruang metrik dan $\{x_n\}_{n=1}^\infty$ barisan dalam $X$.  Maka
$\{ x_n \}_{n=1}^\infty$ konvergen ke $p \in X$ jika dan hanya jika untuk setiap lingkungan terbuka
$U$ dari $p$, terdapat $M \in \N$ sedemikian sehingga untuk semua $n \geq M$,
berlaku $x_n \in U$.
\end{prop}

\begin{proof}
Misalkan $\{ x_n \}_{n=1}^\infty$ konvergen ke $p$.  Misalkan $U$ lingkungan terbuka
dari $p$.  Terdapat $\epsilon > 0$ sedemikian sehingga $B(p,\epsilon) \subset
U$.  Karena barisan tersebut konvergen, pilih $M \in \N$ sedemikian sehingga untuk semua $n \geq
M$, berlaku $d(p,x_n) < \epsilon$, atau dengan kata lain $x_n \in B(p,\epsilon)
\subset U$.

Mari kita buktikan arah sebaliknya.  Untuk $\epsilon > 0$, tetapkan $U \coloneqq
B(p,\epsilon)$ sebagai lingkungan dari $p$.  Maka terdapat $M \in \N$
sedemikian sehingga untuk $n \geq M$, berlaku $x_n \in U = B(p,\epsilon)$, atau dengan kata
lain, $d(p,x_n) < \epsilon$.
\end{proof}

Himpunan tertutup memuat limit setiap barisan konvergen di dalamnya.

\begin{prop} \label{prop:msclosedlim}
Misalkan $(X,d)$ ruang metrik, $E \subset X$ himpunan tertutup,
dan $\{ x_n \}_{n=1}^\infty$ barisan dalam $E$ yang konvergen ke suatu $p \in X$.
Maka $p \in E$.
\end{prop}

\begin{proof}
Mari kita buktikan kontraposisinya.
Misalkan $\{ x_n \}_{n=1}^\infty$ barisan dalam $X$ yang konvergen ke $p \in E^c$.
Karena $E^c$ terbuka, \propref{prop:msconvtopo} menyatakan bahwa terdapat
$M$ sedemikian sehingga untuk semua $n \geq M$,
$x_n \in E^c$.  Jadi, $\{ x_n \}_{n=1}^\infty$ bukan barisan dalam $E$.
\end{proof}

Untuk memperoleh penutupan suatu himpunan $A$, kita mulai dengan $A$, lalu menambahkan 
titik-titik yang merupakan limit barisan dalam $A$.

\begin{prop} \label{prop:msclosureapprseq}
Misalkan $(X,d)$ ruang metrik dan $A \subset X$.
Maka $p \in \widebar{A}$ jika dan hanya jika terdapat barisan
$\{ x_n \}_{n=1}^\infty$ yang terdiri atas
elemen-elemen dalam $A$ sedemikian sehingga $\lim\limits_{n\to\infty} x_n = p$.
\end{prop}

\begin{proof}
Misalkan $p \in \widebar{A}$.  Untuk setiap $n \in \N$,
\propref{prop:msclosureappr} menjamin adanya
titik $x_n \in B(p,\nicefrac{1}{n}) \cap A$.
Karena $d(p,x_n) < \nicefrac{1}{n}$, kita memperoleh $\lim_{n\to\infty} x_n = p$.

Untuk arah sebaliknya, lihat \exerciseref{exercise:reverseclosedseq}.
\end{proof}

\subsection{Latihan}

\begin{exercise} \label{exercise:reverseclosedseq}
Selesaikan pembuktian 
\propref{prop:msclosureapprseq}:
Misalkan $(X,d)$ ruang metrik dan
$A \subset X$.  Misalkan $p \in X$
sedemikian sehingga terdapat barisan $\{ x_n \}_{n=1}^\infty$ dalam $A$ yang konvergen ke $p$.
Buktikan bahwa $p \in \widebar{A}$.
\end{exercise}

\begin{exercise}
\leavevmode
\begin{enumerate}[a)]
\item
Tunjukkan bahwa $d(x,y) \coloneqq \min \bigl\{ 1, \sabs{x-y} \bigr\}$ mendefinisikan metrik pada $\R$.
\item
Tunjukkan bahwa suatu barisan konvergen dalam $(\R,d)$ jika dan hanya jika barisan tersebut konvergen
dalam metrik standar.
\item
Temukan barisan terbatas dalam $(\R,d)$ yang
tidak memiliki subbarisan konvergen.
\end{enumerate}
\end{exercise}

\begin{exercise}
Buktikan \propref{prop:msconvbound}.
\end{exercise}

\begin{exercise}
Buktikan \propref{prop:msconvifa}.
\end{exercise}

\begin{exercise}
Misalkan $\{x_n\}_{n=1}^\infty$ konvergen ke $p$.  Misalkan $f \colon \N
\to \N$ fungsi satu-ke-satu.  Tunjukkan bahwa
$\{ x_{f(n)} \}_{n=1}^\infty$ konvergen ke $p$.
\end{exercise}

\begin{exercise}
Misalkan $(X,d)$ ruang metrik dengan $d$ metrik diskret.  Misalkan 
$\{ x_n \}_{n=1}^\infty$ barisan konvergen dalam $X$.  Tunjukkan bahwa terdapat
$K \in \N$ sedemikian sehingga untuk semua $n \geq K$, berlaku $x_n = x_K$.
\end{exercise}

\begin{exercise}
Himpunan $S \subset X$ disebut \emph{\myindex{rapat}} dalam $X$ jika
$X \subset \widebar{S}$, atau dengan kata lain, jika untuk setiap $p \in X$,
terdapat barisan $\{ x_n \}_{n=1}^\infty$ dalam $S$ yang konvergen ke $p$.
Buktikan bahwa $\R^n$ memuat subhimpunan rapat yang terhitung.
\end{exercise}

\begin{exercise}[Menantang]
Misalkan $\{ U_n \}_{n=1}^\infty$ barisan menurun ($U_{n+1} \subset U_n$ untuk
semua $n$) dari himpunan-himpunan terbuka dalam ruang metrik $(X,d)$ sedemikian sehingga
$\bigcap_{n=1}^\infty U_n = \{ p \}$ untuk suatu $p \in X$.  Misalkan 
$\{ x_n \}_{n=1}^\infty$ barisan titik dalam $X$ sedemikian sehingga $x_n \in U_n$.  Apakah
$\{ x_n \}_{n=1}^\infty$ harus konvergen ke $p$?  Buktikan atau bangun contoh penyangkal.
\end{exercise}

\begin{exercise}
Misalkan $E \subset X$ tertutup dan
$\{ x_n \}_{n=1}^\infty$ barisan dalam $X$ yang konvergen ke $p \in X$.  Andaikan
$x_n \in E$ untuk tak hingga banyaknya nilai $n \in \N$.  Tunjukkan bahwa $p \in E$.
\end{exercise}

\begin{exercise} \label{exercise:extendedrealsmetric}
Ambil $\R^* = \{ -\infty \} \cup \R \cup \{ \infty \}$ sebagai himpunan bilangan real diperluas.
Definisikan
$d(x,y) \coloneqq \babs{ \frac{x}{1+\sabs{x}} - \frac{y}{1+\sabs{y}} }$
untuk $x, y \in \R$, definisikan
$d(\infty,x) \coloneqq \babs{ 1 - \frac{x}{1+\sabs{x}} }$,
$d(-\infty,x) \coloneqq \babs{ 1 + \frac{x}{1+\sabs{x}} }$
untuk semua $x \in \R$.  Tetapkan pula $d(x,\infty) \coloneqq d(\infty,x)$ dan
$d(x,-\infty) \coloneqq d(-\infty,x)$ untuk semua $x \in \R$, serta tetapkan
$d(\infty,-\infty)=d(-\infty,\infty) \coloneqq 2$ dan $d(\infty,\infty)=d(-\infty,-\infty) \coloneqq 0$.
\begin{enumerate}[a)]
\item
Tunjukkan bahwa $(\R^*,d)$ ruang metrik.
\item
Misalkan $\{ x_n \}_{n=1}^\infty$ barisan bilangan real sedemikian sehingga
untuk setiap $M \in \R$, terdapat $N$ sedemikian sehingga
$x_n \geq M$ untuk semua $n \geq N$.  Tunjukkan bahwa $\lim\limits_{n\to\infty} x_n = \infty$ dalam
$(\R^*,d)$.
\item
Tunjukkan bahwa suatu barisan bilangan real konvergen ke bilangan real
dalam $(\R^*,d)$ jika dan
hanya jika barisan tersebut konvergen dalam $\R$ dengan metrik standar.
\end{enumerate}
\end{exercise}

\begin{exercise}
Misalkan $\{ V_n \}_{n=1}^\infty$ barisan himpunan terbuka
dalam $(X,d)$
sedemikian sehingga $V_{n+1} \supset V_n$ untuk semua $n$.  Misalkan $\{ x_n \}_{n=1}^\infty$ barisan
sedemikian sehingga $x_n \in V_{n+1} \setminus V_n$ dan andaikan 
$\{ x_n \}_{n=1}^\infty$ konvergen ke $p \in X$.  Tunjukkan bahwa $p \in \partial V$
dengan $V = \bigcup_{n=1}^\infty V_n$.
\end{exercise}

\begin{exercise}
Buktikan \propref{prop:mssubseq}.
\end{exercise}

\begin{exercise}
Misalkan $(X,d)$ ruang metrik dan $\{ x_n \}_{n=1}^\infty$ barisan dalam $X$.
Buktikan bahwa $\{ x_n \}_{n=1}^\infty$ konvergen ke $p \in X$
jika dan hanya jika
setiap subbarisan dari $\{ x_n \}_{n=1}^\infty$ memiliki subbarisan yang
konvergen ke $p$.
\end{exercise}

\begin{exercise}
Tinjau $\R^n$, dan misalkan $d$ metrik Euklides standar.
Tetapkan $d'(x,y) \coloneqq \sum_{\ell=1}^n \sabs{x_\ell-y_\ell}$ dan
$d''(x,y) \coloneqq \max \bigl\{
\sabs{x_1-y_1},\sabs{x_2-y_2},\cdots,\sabs{x_n-y_n}
\bigr\}$.
\begin{enumerate}[a)]
\item
Gunakan \exerciseref{exercise:mscross} untuk menunjukkan bahwa
$(\R^n,d')$ dan
$(\R^n,d'')$ merupakan ruang metrik.
\item
Misalkan $\{ x_k \}_{k=1}^\infty$ barisan dalam $\R^n$ dan $p \in \R^n$.
Buktikan bahwa pernyataan-pernyataan berikut ekuivalen:
\begin{enumerate}[(1)]
\item
$\{ x_k \}_{k=1}^\infty$ konvergen ke $p$ dalam $(\R^n,d)$.
\item
$\{ x_k \}_{k=1}^\infty$ konvergen ke $p$ dalam $(\R^n,d')$.
\item
$\{ x_k \}_{k=1}^\infty$ konvergen ke $p$ dalam $(\R^n,d'')$.
\end{enumerate}
\end{enumerate}
\end{exercise}

%%%%%%%%%%%%%%%%%%%%%%%%%%%%%%%%%%%%%%%%%%%%%%%%%%%%%%%%%%%%%%%%%%%%%%%%%%%%%%

\sectionnewpage
\section{Kelengkapan dan kekompakan}
\label{sec:metcompact}

%mbxINTROSUBSECTION

\sectionnotes{2 kuliah}

\subsection{Barisan Cauchy dan kelengkapan}

Seperti halnya untuk barisan bilangan real, kita mendefinisikan barisan Cauchy.

\begin{defn}
Misalkan $(X,d)$ ruang metrik.
Barisan $\{ x_n \}_{n=1}^\infty$ dalam $X$ disebut \emph{\myindex{barisan Cauchy}} jika
untuk setiap $\epsilon > 0$, terdapat $M \in \N$ sedemikian sehingga
untuk semua $n \geq M$ dan semua $k \geq M$, berlaku
\begin{equation*}
d(x_n, x_k) < \epsilon .
\end{equation*}
\end{defn}

Definisi tersebut sekali lagi hanyalah pemindahan konsep
dari bilangan real ke ruang metrik.  Barisan bilangan
real merupakan barisan Cauchy dalam pengertian \chapterref{seq:chapter} jika dan hanya jika
barisan tersebut merupakan barisan Cauchy dalam pengertian di atas, asalkan bilangan real dilengkapi dengan
metrik standar $d(x,y) = \sabs{x-y}$.

\begin{prop}
Barisan konvergen dalam ruang metrik merupakan barisan Cauchy.
\end{prop}

\begin{proof}
Misalkan $\{ x_n \}_{n=1}^\infty$ konvergen ke $p$.
Untuk $\epsilon > 0$, terdapat $M$ sedemikian sehingga untuk semua $n \geq M$,
berlaku $d(p,x_n) < \nicefrac{\epsilon}{2}$.  Jadi,
untuk semua $n,k \geq M$, berlaku
$d(x_n,x_k) \leq d(x_n,p) + d(p,x_k) < \nicefrac{\epsilon}{2} +
\nicefrac{\epsilon}{2} = \epsilon$.
\end{proof}

\begin{defn}
Ruang metrik $(X,d)$ disebut
\emph{\myindex{lengkap}} atau \emph{\myindex{lengkap secara Cauchy}}
jika setiap barisan Cauchy $\{ x_n \}_{n=1}^\infty$ dalam $X$
konvergen ke suatu $p \in X$.
\end{defn}

\begin{prop}
Ruang $\R^n$ dengan metrik standar merupakan ruang metrik lengkap.
\end{prop}

Untuk $\R = \R^1$, kelengkapan telah dibuktikan dalam \chapterref{seq:chapter}.
Pembuktian kelengkapan dalam $\R^n$
merupakan reduksi ke kasus satu dimensi.

\begin{proof}
Misalkan $\{ x_m \}_{m=1}^\infty$ barisan Cauchy
dalam $\R^n$, dengan $x_m = \bigl(x_{m,1},x_{m,2},\ldots,x_{m,n}\bigr) \in \R^n$.
Karena barisan tersebut merupakan barisan Cauchy, untuk $\epsilon > 0$ terdapat $M$ sedemikian sehingga untuk semua
$i,j \geq M$,
\begin{equation*}
d(x_i,x_j) < \epsilon.
\end{equation*}

Tetapkan sebarang $k=1,2,\ldots,n$.  Untuk $i,j \geq M$,
\begin{equation*}
\bigl\lvert x_{i,k} - x_{j,k} \bigr\rvert
=
\sqrt{{\bigl(x_{i,k} - x_{j,k}\bigr)}^2}
\leq
\sqrt{\sum_{\ell=1}^n {\bigl(x_{i,\ell}-x_{j,\ell}\bigr)}^2}
= d(x_i,x_j) < \epsilon .
\end{equation*}
Jadi, barisan $\{ x_{m,k} \}_{m=1}^\infty$ merupakan barisan Cauchy.  Karena $\R$
lengkap, barisan tersebut konvergen; terdapat $y_k \in \R$ sedemikian sehingga
$y_k = \lim_{m\to\infty} x_{m,k}$.
Tuliskan $y = (y_1,y_2,\ldots,y_n) \in \R^n$.
Menurut \propref{prop:msconveuc}, $\{ x_m \}_{m=1}^\infty$ konvergen
ke $y \in \R^n$, sehingga $\R^n$ lengkap.
\end{proof}

Dalam bahasa ruang metrik,
hasil tentang kekontinuan dalam \sectionref{sec:liminter}
menyatakan bahwa ruang metrik
$C\bigl([a,b],\R\bigr)$ dari \exampleref{example:msC01} lengkap.
Pembuktiannya diperoleh dengan \myquote{menguraikan definisi-definisi} 
dan diserahkan sebagai \exerciseref{exercise:CabRcomplete}.

\begin{prop} \label{prop:CabRcomplete}
Ruang fungsi kontinu $C\bigl([a,b],\R\bigr)$ dengan metrik yang diinduksi oleh norma
seragam merupakan ruang metrik lengkap.
\end{prop}

Subhimpunan dari ruang metrik lengkap seperti $\R^n$, dengan metrik subruang, belum tentu
lengkap.  Sebagai contoh, $(0,1]$ dengan metrik subruang tidak
lengkap karena $\{ \nicefrac{1}{n} \}_{n=1}^\infty$ merupakan barisan Cauchy dalam $(0,1]$
yang tidak memiliki limit dalam $(0,1]$.
Namun, 
subruang tertutup dari ruang metrik lengkap bersifat
lengkap.  Salah satu cara
memahami himpunan tertutup adalah bahwa himpunan tersebut memuat semua titik yang
dapat dicapai dari dalam himpunan melalui suatu barisan.
Pembuktiannya diserahkan sebagai 
\exerciseref{exercise:closedcomplete}.

\begin{prop} \label{prop:closedcomplete}
Misalkan $(X,d)$ ruang metrik lengkap dan $E \subset X$
tertutup.
Maka $E$ merupakan ruang metrik lengkap dengan metrik subruang.
\end{prop}

\subsection{Kekompakan}

\begin{defn}
Misalkan $(X,d)$ ruang metrik dan $K \subset X$. 
Himpunan $K$ disebut \emph{\myindex{kompak}}
jika untuk setiap koleksi
himpunan terbuka $\{ U_{\lambda} \}_{\lambda \in I}$ sedemikian sehingga
\begin{equation*}
K \subset \bigcup_{\lambda \in I} U_\lambda ,
\end{equation*}
terdapat himpunan bagian berhingga
$\{ \lambda_1, \lambda_2,\ldots,\lambda_m \} \subset I$
sedemikian sehingga
\begin{equation*}
K \subset \bigcup_{j=1}^m U_{\lambda_j} .
\end{equation*}
\end{defn}

Koleksi himpunan terbuka $\{ U_{\lambda} \}_{\lambda \in I}$ seperti di atas
disebut \emph{\myindex{selimut terbuka}} bagi $K$.  Cara lain untuk mengatakan bahwa
$K$ kompak ialah dengan mengatakan bahwa \emph{setiap selimut terbuka bagi $K$ memiliki
\myindex{subselimut} berhingga}.

\begin{example}
Misalkan $\R$ ruang metrik dengan metrik standar.

Himpunan $\R$ tidak kompak.  Bukti: Untuk $j \in \N$, misalkan $U_j \coloneqq (-j,j)$.
Setiap $x \in \R$ berada dalam suatu $U_j$ (berdasarkan
\hyperref[thm:arch:i]{sifat Archimedes}), sehingga kita memperoleh sebuah selimut terbuka.
Andaikan kita memiliki
subselimut berhingga $\R \subset U_{j_1} \cup U_{j_2} \cup \cdots \cup U_{j_m}$,
dan andaikan $j_1 < j_2 < \cdots < j_m$.  Maka $\R \subset U_{j_m}$, tetapi ini
kontradiksi karena $j_m \in \R$ di satu pihak dan $j_m \notin U_{j_m} =
(-j_m,j_m)$ di pihak
lain.

Himpunan $(0,1) \subset \R$ juga tidak kompak.  Bukti:  Ambil 
himpunan-himpunan $U_{j} \coloneqq (\nicefrac{1}{j},1-\nicefrac{1}{j})$ untuk $j=3,4,5,\ldots$.
Seperti di atas, $(0,1) = \bigcup_{j=3}^\infty U_j$.  Dan seperti di atas pula,
jika terdapat subselimut berhingga, maka terdapat satu $U_j$ sedemikian sehingga $(0,1)
\subset U_j$, yang kembali menghasilkan kontradiksi.

Himpunan $\{ 0 \} \subset \R$ kompak.  Bukti: Diberikan selimut terbuka $\{
U_{\lambda} \}_{\lambda \in I}$, harus terdapat $\lambda_0 \in I$ sedemikian sehingga $0
\in U_{\lambda_0}$
karena $\{ U_{\lambda} \}_{\lambda \in I}$ merupakan selimut.
Hanya dengan $U_{\lambda_0}$, kita memperoleh
subselimut berhingga.

Kita akan membuktikan di bawah bahwa $[0,1]$, bahkan setiap interval tertutup
dan terbatas $[a,b]$, bersifat kompak.
\end{example}

\begin{prop}
Misalkan $(X,d)$ ruang metrik.  Jika $K \subset X$ kompak, maka
$K$ tertutup dan terbatas.
\end{prop}

\begin{proof}
Pertama, kita buktikan bahwa himpunan kompak bersifat terbatas.
Tetapkan $p \in X$.  Kita memiliki selimut terbuka
\begin{equation*}
K \subset \bigcup_{n=1}^\infty B(p,n) = X .
\end{equation*}
Jika $K$ kompak, maka terdapat sejumlah indeks
$n_1 < n_2 < \ldots < n_m$ sedemikian sehingga
\begin{equation*}
K \subset \bigcup_{j=1}^m B(p,n_j) = B(p,n_m) .
\end{equation*}
Karena $K$ termuat dalam sebuah bola, $K$ terbatas.
Lihat bagian kiri \figureref{fig:compactbndclosed}.

Selanjutnya, kita tunjukkan bahwa himpunan yang tidak tertutup tidak kompak.  Andaikan 
$\widebar{K} \neq K$, yaitu terdapat titik $x \in \widebar{K}
\setminus K$.
Jika $y \neq x$, maka
$y \notin C(x,\nicefrac{1}{n})$
untuk $n \in \N$
sedemikian sehingga $\nicefrac{1}{n} < d(x,y)$.
Lebih lanjut, $x \notin K$, sehingga
\begin{equation*}
K \subset \bigcup_{n=1}^\infty {C(x,\nicefrac{1}{n})}^c .
\end{equation*}
Bola tertutup bersifat tertutup, sehingga komplemennya ${C(x,\nicefrac{1}{n})}^c$ terbuka, dan
kita memperoleh sebuah selimut terbuka.
Jika kita mengambil sebarang
koleksi berhingga indeks $n_1 < n_2 < \ldots < n_m$, maka 
\begin{equation*}
\bigcup_{j=1}^m {C(x,\nicefrac{1}{n_j})}^c 
=
{C(x,\nicefrac{1}{n_m})}^c 
\end{equation*}
Karena $x$ berada dalam penutupan $K$,
maka
$C(x,\nicefrac{1}{n_m}) \cap K \neq \emptyset$.  Jadi, tidak ada
subselimut berhingga dan $K$ tidak kompak.
Lihat bagian kanan \figureref{fig:compactbndclosed}.
\end{proof}

\begin{myfigureht}
\myincludepdft{compact_bnd_closed}{%
Dua diagram.  Himpunan K digambarkan sebagai daerah berarsir dengan
batas bergaris utuh.  Pada diagram kiri, titik p ditonjolkan
bersama 3 bola terbuka bersarang (dalam hal ini berbentuk lingkaran)
berjari-jari 1, 2, dan 3, semuanya berpusat di p dan berbatas garis putus-putus.
Himpunan K terletak di dalam
bola terbesar.
Pada diagram kanan, latar belakang diagram diarsir.
Titik x yang berada pada batas K diperlihatkan, dan 4 bola tertutup
digambar sebagai lingkaran dengan bagian dalam yang diarsir dalam warna abu-abu
yang semakin muda, menandakan bahwa yang sedang ditinjau adalah komplemennya.
Bola terluar berjari-jari 1, lalu terdapat bola berjari-jari setengah,
kemudian bola berjari-jari sepertiga, lalu bola berjari-jari seperempat.
Kita perhatikan bahwa karena x terletak pada batas, komplemen bola-bola tersebut
menutupi bagian K yang semakin besar.  Di sisi lain, kita juga perhatikan bahwa K
beririsan tak kosong dengan setiap bola tertutup tersebut.}
\caption{Pembuktian bahwa himpunan kompak bersifat terbatas (kiri) dan tertutup (kanan).\label{fig:compactbndclosed}}
\end{myfigureht}

Kita buktikan di bawah bahwa 
dalam ruang Euklides berdimensi hingga,
setiap himpunan tertutup dan terbatas bersifat kompak.
Jadi, himpunan-himpunan tertutup dan terbatas
dalam $\R^n$ merupakan contoh himpunan kompak.
Tidak benar bahwa dalam setiap ruang metrik, sifat tertutup dan terbatas ekuivalen
dengan kekompakan.  Contoh sederhana ialah ruang metrik tak lengkap seperti
$(0,1)$ dengan metrik subruang yang berasal dari $\R$.
Namun, pada banyak ruang metrik lengkap yang sangat berguna,
sifat tertutup dan terbatas tidak
cukup untuk menjamin kekompakan: $C\bigl([a,b],\R\bigr)$ merupakan ruang metrik
lengkap, tetapi bola satuan tertutup $C(0,1)$ tidak kompak; lihat
\exerciseref{exercise:msclbounnotcompt}.  Bandingkan pula dengan
\exerciseref{exercise:mstotbound}.

Dengan mengesampingkan sifat terbatas untuk sementara,
perhatikan pula perbedaan antara sifat tertutup dan sifat kompak.  Sifat
tertutup bergantung pada ruang metrik ambien:
Himpunan $(0,1]$ tidak tertutup dalam $\R$,
tetapi tertutup dalam subruang $(0,\infty)$.
Namun, suatu himpunan $K$ kompak
dalam suatu ruang metrik $(X,d)$ jika dan hanya jika ia kompak dalam metrik
subruang pada $K$.  Jadi, untuk himpunan kompak, kita tidak perlu menanyakan
ruang metrik mana yang memuatnya.
Di sisi lain, setiap himpunan selalu tertutup dalam 
metrik subruang ketika dipandang sebagai himpunan bagian dari dirinya sendiri.
Lihat juga \exerciseref{exercise:subspacecompact}.

Sifat yang berguna dari himpunan kompak dalam ruang metrik ialah bahwa setiap
barisan dalam himpunan tersebut memiliki subbarisan konvergen yang konvergen
ke suatu titik dalam himpunan tersebut.
Himpunan seperti ini disebut
\emph{\myindex{kompak secara sekuensial}}.
Kita akan membuktikan bahwa dalam
konteks ruang metrik, suatu himpunan kompak jika dan hanya jika kompak secara
sekuensial.
Pertama, kita buktikan sebuah lema.

\begin{lemma}[Lema selimut Lebesgue%
\footnote{Dinamai menurut matematikawan Prancis
\href{https://en.wikipedia.org/wiki/Henri_Lebesgue}{Henri L\'eon Lebesgue}
(1875--1941).
Bilangan $\delta$ kadang-kadang disebut \myindex{bilangan Lebesgue} dari
selimut tersebut.}]\label{ms:lebesgue}
\index{Lema selimut Lebesgue}
Misalkan $(X,d)$ ruang metrik dan $K \subset X$.  Andaikan 
setiap barisan dalam $K$ memiliki subbarisan yang konvergen dalam $K$.  Diberikan
selimut terbuka $\{ U_\lambda \}_{\lambda \in I}$ bagi $K$, terdapat
$\delta > 0$ sedemikian sehingga untuk setiap $x \in K$, terdapat $\lambda \in I$
dengan $B(x,\delta) \subset U_\lambda$.
\end{lemma}

\begin{proof}
Kita buktikan lema ini melalui kontrapositifnya.
Jika kesimpulannya tidak benar, maka
terdapat
selimut terbuka $\{ U_\lambda \}_{\lambda \in I}$ bagi $K$ dengan
sifat berikut.
Untuk setiap $n \in \N$, terdapat $x_n \in K$ sedemikian sehingga
$B(x_n,\nicefrac{1}{n})$ tidak termuat dalam satu pun $U_\lambda$.
Ambil sebarang $x \in K$.  Terdapat
$\lambda \in I$ sedemikian sehingga $x \in U_\lambda$.  Karena $U_\lambda$ terbuka,
terdapat $\epsilon > 0$ 
sedemikian sehingga $B(x,\epsilon) \subset U_\lambda$.  Ambil $M$ sedemikian sehingga
$\nicefrac{1}{M} < \nicefrac{\epsilon}{2}$.  Jika $y \in 
B(x,\nicefrac{\epsilon}{2})$ dan $n \geq M$, maka 
\begin{equation*}
B(y,\nicefrac{1}{n}) \subset
B(y,\nicefrac{1}{M}) \subset
B(y,\nicefrac{\epsilon}{2}) \subset B(x,\epsilon)
\subset U_\lambda ,
\end{equation*}
di mana 
$B(y,\nicefrac{\epsilon}{2}) \subset B(x,\epsilon)$
merupakan akibat ketaksamaan segitiga.
Lihat \figureref{fig:lebesguedelta}.
Jadi, $y \neq x_n$.
Dengan kata lain, untuk semua $n \geq M$, berlaku $x_n \notin B(x,\nicefrac{\epsilon}{2})$. 
Barisan tersebut tidak dapat memiliki subbarisan yang konvergen ke $x$.  Karena $x \in K$
dipilih sebarang, pembuktian selesai.
\end{proof}

\begin{myfigureht}
\myincludepdft{lebesguedelta}{%
Diagram sebuah himpunan terbuka berarsir U subskrip lambda.
Titik x dalam U subskrip lambda ditandai, dan sebuah bola berjari-jari
epsilon yang berpusat di x digambar; bola ini terletak di dalam U subskrip lambda.
Bola berjari-jari epsilon per 2 yang berpusat di x juga digambar dan
terletak di dalam bola yang lebih besar.  Titik y dalam bola yang lebih kecil ditandai,
dan sebuah bola berjari-jari epsilon per 2 yang berpusat di y diperlihatkan.  Kita perhatikan
bahwa bola yang berpusat di y ini seluruhnya terletak di dalam bola
berjari-jari epsilon yang berpusat di x, yaitu bola terbesar, dan karena itu juga di dalam U
subskrip lambda.  Terakhir, diperlihatkan sebuah bola kecil berjari-jari 1 per n yang berpusat di y
dan terletak di dalam bola epsilon per 2 yang berpusat di y.}
\caption{Pembuktian lema selimut Lebesgue.
Perhatikan bahwa $B(y,\nicefrac{\epsilon}{2}) \subset
B(x,\epsilon)$ berdasarkan ketaksamaan segitiga.\label{fig:lebesguedelta}}
\end{myfigureht}

Penting untuk memahami apa yang dinyatakan oleh lema tersebut.  Lema itu menyatakan bahwa
jika $K$ kompak secara sekuensial, maka untuk setiap
selimut yang diberikan, terdapat satu $\delta > 0$.  Nilai $\delta$ bergantung pada selimutnya,
tetapi, tentu saja, tidak bergantung pada $x$.

Sebagai contoh, misalkan $K \coloneqq [-10,10]$ dan misalkan
$U_n \coloneqq (n,n+2)$ untuk $n \in \Z$
membentuk sebuah selimut terbuka.
Tinjau $x \in K$. Terdapat $n \in \Z$ 
sedemikian sehingga $n \leq x < n+1$.
Jika $n \leq x < n+\nicefrac{1}{2}$, maka
$B\bigl(x,\nicefrac{1}{2}\bigr) \subset U_{n-1}$.
Jika $n+ \nicefrac{1}{2} \leq x < n+1$, maka
$B\bigl(x,\nicefrac{1}{2}\bigr) \subset U_{n}$.  Jadi, $\delta =
\nicefrac{1}{2}$ dapat digunakan.  Himpunan-himpunan 
$U'_n \coloneqq \bigl(\frac{n}{2},\frac{n+2}{2} \bigr)$,
kembali membentuk sebuah selimut terbuka,
tetapi kini nilai $\delta$ terbesar yang dapat digunakan adalah $\nicefrac{1}{4}$.

Di sisi lain, $\N \subset \R$ tidak kompak secara sekuensial.
Sebagai latihan, carilah selimut yang tidak memungkinkan sebarang $\delta > 0$.


\begin{thm} \label{thm:mscompactisseqcpt}
Misalkan $(X,d)$ ruang metrik.  Maka $K \subset X$ kompak jika dan
hanya jika setiap barisan dalam $K$ memiliki subbarisan yang konvergen ke
suatu titik dalam $K$.
\end{thm}

\begin{proof}
Klaim: \emph{Misalkan $K \subset X$ himpunan bagian dari $X$ dan
$\{ x_n \}_{n=1}^\infty$ barisan dalam $K$.  Andaikan bahwa untuk setiap $x \in K$,
terdapat bola $B(x,\alpha_x)$ untuk suatu $\alpha_x > 0$ sedemikian sehingga
$x_n \in B(x,\alpha_x)$ hanya untuk berhingga banyak $n \in \N$.
Maka $K$ tidak kompak.}

Pembuktian klaim:
Perhatikan bahwa
\begin{equation*}
K \subset \bigcup_{x \in K} B(x,\alpha_x) .
\end{equation*}
Setiap koleksi berhingga dari bola-bola tersebut hanya memuat berhingga banyak
suku dari $\{ x_n \}_{n=1}^\infty$,
sehingga pasti terdapat $x_n \in K$
yang tidak berada dalam gabungannya.  Jadi, $K$ tidak kompak dan klaim
terbukti.

Sekarang andaikan bahwa $K$ kompak dan $\{ x_n \}_{n=1}^\infty$ barisan dalam $K$.
Maka terdapat $x \in K$ sedemikian sehingga
untuk setiap $\delta > 0$,
$B(x,\delta)$ memuat $x_n$ untuk tak berhingga banyak $n \in \N$.
Kita definisikan subbarisan tersebut secara induktif.
Bola $B(x,1)$ memuat suatu $x_k$, jadi tetapkan $n_1 \coloneqq k$.
Andaikan $n_{j-1}$ telah didefinisikan.
Pasti terdapat $k > n_{j-1}$
sedemikian sehingga $x_k \in B(x,\nicefrac{1}{j})$.  Definisikan
$n_j \coloneqq k$.
Kini kita memperoleh subbarisan $\{ x_{n_j} \}_{j=1}^\infty$.
Karena
$d(x,x_{n_j}) < \nicefrac{1}{j}$,  \propref{prop:msconvifa} menyatakan bahwa
$\lim_{j\to\infty} x_{n_j} = x$.

Untuk arah sebaliknya, andaikan setiap barisan dalam~$K$
memiliki 
subbarisan yang konvergen dalam~$K$.
Ambil
selimut terbuka $\{ U_\lambda \}_{\lambda \in I}$ bagi $K$.
Dengan menggunakan lema selimut Lebesgue di atas, diperoleh $\delta > 0$
sedemikian sehingga untuk setiap $x \in K$, terdapat $\lambda \in I$ dengan
$B(x,\delta) \subset U_\lambda$.

Pilih $x_1 \in K$ dan dapatkan $\lambda_1 \in I$ sedemikian sehingga $B(x_1,\delta) \subset
U_{\lambda_1}$.
Jika $K \subset U_{\lambda_1}$, kita berhenti karena telah memperoleh
subselimut berhingga.
Jika tidak, pasti terdapat
titik $x_2 \in K \setminus U_{\lambda_1}$.
Perhatikan bahwa $d(x_2,x_1) \geq \delta$.
Pasti terdapat $\lambda_2 \in I$ sedemikian sehingga
$B(x_2,\delta) \subset U_{\lambda_2}$.
Kita lanjutkan secara induktif.  Andaikan $\lambda_{n-1}$ telah didefinisikan.
Entah
$U_{\lambda_1} \cup
U_{\lambda_2} \cup \cdots \cup
U_{\lambda_{n-1}}$ merupakan selimut berhingga bagi $K$, sehingga kita
berhenti, atau
pasti terdapat 
titik $x_n \in K \setminus \bigl( U_{\lambda_1} \cup
U_{\lambda_2} \cup \cdots \cup
U_{\lambda_{n-1}}\bigr)$.
Perhatikan bahwa $d(x_n,x_j) \geq \delta$ untuk semua $j = 1,2,\ldots,n-1$.
Selanjutnya, pasti terdapat $\lambda_n \in I$
sedemikian sehingga $B(x_n,\delta) \subset U_{\lambda_n}$.
Lihat \figureref{fig:seqcompactiscompact}.

\begin{myfigureht}
\myincludepdft{seqcompactiscompact}{%
Sebuah daerah berarsir dengan batas berupa garis penuh diberi label K.
Empat titik di K diberi label x subskrip 1, x subskrip 2, x subskrip 3, dan x subskrip 4.
Lingkaran berjari-jari delta digambar dengan pusat di x subskrip 1, x subskrip 2, dan x subskrip 3.
Di sekeliling setiap lingkaran tersebut terdapat tiga daerah dengan batas berupa garis
putus-putus yang diberi label U subskrip lambda subskrip 1, 
U subskrip lambda subskrip 2, dan
U subskrip lambda subskrip 3.
Cakram-cakram tersebut, dan karena itu daerah-daerah yang mengelilinginya, tampak secara bertahap
menyelimuti seluruh himpunan K dari sisi kiri gambar.}
\caption{Menyelimuti $K$ dengan $U_{\lambda}$.  Titik-titik
$x_1,x_2,x_3,x_4$, 
ketiga himpunan 
$U_{\lambda_1}$,
$U_{\lambda_2}$,
$U_{\lambda_3}$,
dan 
ketiga bola pertama
yang berjari-jari $\delta$ digambar.\label{fig:seqcompactiscompact}}
\end{myfigureht}

Pada suatu tahap kita memperoleh subselimut berhingga bagi $K$,
atau kita memperoleh
barisan tak berhingga
$\{ x_n \}_{n=1}^\infty$ seperti di atas.
Untuk memperoleh kontradiksi, andaikan bahwa
tidak terdapat subselimut berhingga dan kita memperoleh barisan $\{ x_n \}_{n=1}^\infty$.
Untuk semua $n$ dan $k$ dengan $n \neq k$, 
berlaku $d(x_n,x_k) \geq \delta$.
Jadi, tidak ada subbarisan dari $\{ x_n \}_{n=1}^\infty$ yang merupakan barisan Cauchy.
Dengan demikian, tidak ada subbarisan dari $\{ x_n \}_{n=1}^\infty$ yang konvergen,
dan ini merupakan kontradiksi.
\end{proof}

\begin{example} \label{example:intervalcompact}
\thmref{thm:bwseq},
yaitu teorema Bolzano--Weierstrass untuk barisan bilangan real,
menyatakan bahwa setiap barisan terbatas dalam $\R$ memiliki
subbarisan yang konvergen.  Oleh karena itu, setiap barisan dalam interval tertutup $[a,b] \subset \R$ memiliki 
subbarisan yang konvergen.  Limitnya juga berada dalam $[a,b]$ karena operasi pengambilan limit
mempertahankan ketaksamaan tak ketat.  Jadi, interval tertutup dan terbatas $[a,b]
\subset \R$ kompak (secara sekuensial).
\end{example}

\begin{prop} \label{prop:closedsubsetcompact}
Misalkan $(X,d)$ ruang metrik dan misalkan $K \subset X$ kompak.  Jika
$E \subset K$ himpunan tertutup, maka $E$ kompak.
\end{prop}

Karena $K$ tertutup, $E$ tertutup dalam $K$ jika dan
hanya jika himpunan tersebut tertutup dalam $X$.
Lihat \propref{prop:topology:subspacesame}.

\begin{proof}
Misalkan $\{ x_n \}_{n=1}^\infty$ barisan dalam $E$.  Barisan ini juga merupakan barisan dalam $K$.
Oleh karena itu, barisan tersebut memiliki subbarisan $\{ x_{n_j} \}_{j=1}^\infty$ yang konvergen ke
suatu $x \in K$.  Karena $E$ tertutup, limit barisan dalam $E$ juga berada dalam $E$
sehingga $x \in E$.  Jadi, $E$ kompak.
\end{proof}

\begin{thm}[Heine--Borel%
\footnote{Dinamai menurut matematikawan Jerman 
\href{https://en.wikipedia.org/wiki/Eduard_Heine}{Heinrich Eduard Heine}
(1821--1881),
dan matematikawan Prancis
\href{https://en.wikipedia.org/wiki/\%C3\%89mile_Borel}{F\'elix \'Edouard Justin \'Emile Borel}
(1871--1956).}]%
\index{teorema Heine--Borel}
\label{thm:msbw}
Setiap himpunan bagian $K \subset \R^n$ yang tertutup dan terbatas bersifat kompak.
\end{thm}

Jadi, himpunan bagian dari $\R^n$ kompak jika dan hanya jika tertutup dan terbatas,
suatu syarat yang jauh lebih mudah diperiksa.
Perlu ditegaskan kembali bahwa teorema Heine--Borel hanya berlaku untuk $\R^n$, bukan
untuk ruang metrik secara umum.
Teorema ini bahkan tidak berlaku untuk subruang dari $\R^n$; teorema tersebut hanya berlaku dalam
$\R^n$ itu sendiri.
Secara umum, sifat kompak mengakibatkan sifat tertutup dan
terbatas, tetapi tidak sebaliknya.

\begin{proof}
Untuk $\R = \R^1$, andaikan $K \subset \R$ tertutup dan terbatas.
Maka $K \subset [a,b]$ untuk suatu interval tertutup dan terbatas,
yang kompak menurut \exampleref{example:intervalcompact}.
Karena $K$ merupakan himpunan bagian tertutup dari suatu himpunan kompak,
maka himpunan ini kompak menurut \propref{prop:closedsubsetcompact}.

Kita buktikan kasus $n=2$ dan serahkan kasus $n$ sebarang sebagai latihan.
Karena $K \subset \R^2$ terbatas, terdapat himpunan
$B=[a,b]\times[c,d] \subset \R^2$ sedemikian sehingga $K \subset B$.  Kita akan menunjukkan
bahwa $B$ kompak.  Karena itu, $K$, sebagai himpunan bagian tertutup dari $B$ yang kompak,
juga kompak.  

Misalkan $\bigl\{ (x_k,y_k) \bigr\}_{k=1}^\infty$ barisan dalam $B$.  Dengan kata lain,
$a \leq x_k \leq b$ dan
$c \leq y_k \leq d$ untuk semua $k$.  Setiap barisan bilangan real yang terbatas
memiliki subbarisan
yang konvergen, sehingga terdapat subbarisan $\{ x_{k_j} \}_{j=1}^\infty$
yang konvergen.  Subbarisan 
$\{ y_{k_j} \}_{j=1}^\infty$ juga merupakan barisan terbatas sehingga terdapat
subbarisan
$\{ y_{k_{j_i}} \}_{i=1}^\infty$ yang konvergen.  Subbarisan dari suatu
barisan konvergen tetap konvergen, sehingga 
$\{ x_{k_{j_i}} \}_{i=1}^\infty$ konvergen.
Tetapkan
\begin{equation*}
x \coloneqq \lim_{i\to\infty} x_{k_{j_i}}
\qquad \text{dan} \qquad
y \coloneqq \lim_{i\to\infty} y_{k_{j_i}} .
\end{equation*}
Menurut \propref{prop:msconveuc},
$\bigl\{ (x_{k_{j_i}},y_{k_{j_i}}) \bigr\}_{i=1}^\infty$ konvergen ke $(x,y)$.
Selain itu, karena $a \leq x_k \leq b$ dan
$c \leq y_k \leq d$ untuk semua $k$, kita tahu bahwa $(x,y) \in B$.
\end{proof}

\begin{example}
Metrik diskret kembali memberikan contoh-contoh tandingan yang menarik.
Misalkan $(X,d)$ ruang metrik dengan metrik diskret, yaitu $d(x,y) = 1$
jika $x \neq y$.  Andaikan
$X$ himpunan tak berhingga.  Maka
\begin{enumerate}[(i)]
\item $(X,d)$ merupakan ruang metrik lengkap.
\item Setiap himpunan bagian $K \subset X$ tertutup dan terbatas.
\item Himpunan bagian $K \subset X$ kompak jika dan hanya jika himpunan tersebut berhingga.
\item Kesimpulan lema selimut Lebesgue selalu dipenuhi,
misalnya dengan $\delta = \nicefrac{1}{2}$,
bahkan untuk $K \subset X$ yang tidak kompak.
\end{enumerate}
Pembuktian
pernyataan-pernyataan di atas bersifat langsung atau diserahkan sebagai latihan
di bawah ini.
\end{example}

\begin{remark}
Satu hal yang perlu diperhatikan mengenai barisan Cauchy, kelengkapan, kekompakan,
dan kekonvergenan ialah bahwa kekompakan dan kekonvergenan hanya bergantung pada
topologi, yaitu pada himpunan mana saja yang merupakan himpunan terbuka.  Di sisi lain,
barisan Cauchy dan kelengkapan bergantung pada metrik yang digunakan.
Lihat \exerciseref{exercise:cauchydepndsonmetric}.
\end{remark}

\subsection{Latihan}

\begin{exercise}
Misalkan $(X,d)$ ruang metrik dan $A$ subhimpunan berhingga dari $X$.
Tunjukkan bahwa $A$ kompak.
\end{exercise}

\begin{samepage}
\begin{exercise}
Misalkan $A \coloneqq \{ \nicefrac{1}{n} : n \in \N \} \subset \R$.
\begin{enumerate}[a)]
\item
Tunjukkan secara langsung dari definisi bahwa $A$
tidak kompak.
\item
Tunjukkan secara langsung dari definisi bahwa $A \cup \{ 0 \}$
kompak.
\end{enumerate}
\end{exercise}
\end{samepage}


\begin{exercise}
Misalkan $(X,d)$ ruang metrik dengan metrik diskret.
\begin{enumerate}[a)]
\item
Buktikan bahwa $X$ lengkap.
\item
Buktikan bahwa $X$ kompak jika dan hanya jika $X$ himpunan berhingga.
\end{enumerate}
\end{exercise}

\begin{exercise}
\leavevmode
\begin{enumerate}[a)]
\item
Tunjukkan bahwa gabungan berhingga dari himpunan-himpunan kompak merupakan himpunan kompak.
\item
Temukan contoh yang menunjukkan bahwa gabungan tak berhingga banyak himpunan kompak tidak selalu
kompak.
\end{enumerate}
\end{exercise}

\begin{exercise}
Buktikan \thmref{thm:msbw} untuk dimensi sebarang.
Petunjuk: Kuncinya adalah menggunakan notasi yang tepat.
\end{exercise}

\begin{exercise} \label{exercise:subspacecompact}
Tunjukkan bahwa himpunan kompak $K$ (dalam sebarang ruang metrik)
merupakan ruang metrik lengkap (dengan metrik subruang).
\end{exercise}

\begin{exercise} \label{exercise:CabRcomplete}
Misalkan $C\bigl([a,b],\R\bigr)$ ruang metrik seperti dalam \exampleref{example:msC01}.  Tunjukkan bahwa
$C\bigl([a,b],\R\bigr)$ merupakan ruang metrik lengkap.
\end{exercise}

\begin{exercise}[Menantang] \label{exercise:msclbounnotcompt}
Misalkan $C\bigl([0,1],\R\bigr)$ ruang metrik pada \exampleref{example:msC01}.  Misalkan $0$
menyatakan fungsi nol.  Tunjukkan bahwa bola tertutup
$C(0,1)$ tidak kompak (meskipun tertutup dan terbatas).
Petunjuk: Bangun barisan fungsi kontinu yang berbeda
$\{ f_n \}_{n=1}^\infty$ sedemikian sehingga
$d(f_n,0) = 1$ dan $d(f_n,f_k) = 1$ untuk semua $n \neq k$.  Tunjukkan bahwa
himpunan $\{ f_n : n \in \N \} \subset C(0,1)$ tertutup tetapi tidak kompak.
Lihat \chapterref{fs:chapter} untuk memperoleh inspirasi.
\end{exercise}

\begin{exercise}[Menantang]
Tunjukkan bahwa terdapat metrik pada $\R$ yang membuat $\R$ menjadi ruang metrik kompak.
\end{exercise}

\begin{exercise}
Misalkan $(X,d)$ lengkap dan terdapat koleksi tak hingga terhitung
dari himpunan kompak tak kosong $E_1 \supset E_2 \supset E_3 \supset
\cdots$.  Buktikan bahwa $\bigcap_{j=1}^\infty E_j \neq \emptyset$.
\end{exercise}

\begin{exercise}[Menantang]
Misalkan $C\bigl([0,1],\R\bigr)$ ruang metrik pada \exampleref{example:msC01}.
Misalkan $K$ himpunan semua $f \in C\bigl([0,1],\R\bigr)$ sedemikian sehingga
$f$ sama dengan polinom kuadrat, yaitu\ $f(x) = a+bx+cx^2$, dan
$\babs{f(x)} \leq 1$ untuk semua $x \in [0,1]$,
yakni $f \in C(0,1)$.  Tunjukkan bahwa $K$ kompak.
\end{exercise}

\begin{exercise}[Menantang] \label{exercise:mstotbound}
Misalkan $(X,d)$ ruang metrik lengkap.
Tunjukkan bahwa $K \subset X$ kompak jika dan hanya jika $K$ tertutup
dan untuk setiap $\epsilon > 0$
terdapat berhingga banyak titik $x_1,x_2,\ldots,x_n$ dengan
$K \subset \bigcup_{j=1}^n B(x_j,\epsilon)$.
Catatan: Himpunan $K$ seperti itu disebut \emph{\myindex{terbatas total}},
sehingga dalam ruang metrik lengkap, suatu himpunan kompak jika dan hanya
jika tertutup dan terbatas total.
\end{exercise}

\begin{exercise}
Ambil $\N \subset \R$ dengan metrik standar.  Temukan selimut terbuka bagi $\N$
sedemikian sehingga kesimpulan lemma selimut Lebesgue tidak berlaku.
\end{exercise}

\begin{exercise}
Buktikan bentuk umum \myindex{teorema Bolzano--Weierstrass}:
Setiap barisan terbatas $\{ x_k \}_{k=1}^\infty$ dalam $\R^n$
memiliki subbarisan konvergen.
\end{exercise}

\begin{exercise}
Misalkan $X$ ruang metrik dan
$C \subset \sP(X)$ adalah himpunan semua subhimpunan kompak tak kosong dari $X$.
Dengan menggunakan metrik Hausdorff dari \exerciseref{exercise:mshausdorffpseudo},
tunjukkan bahwa $(C,d_H)$ merupakan ruang metrik.  Artinya, tunjukkan bahwa
jika $L$ dan $K$ subhimpunan kompak tak kosong, maka $d_H(L,K) = 0$
jika dan hanya jika $L=K$.
\end{exercise}

\begin{exercise} \label{exercise:closedcomplete}
Buktikan \propref{prop:closedcomplete}.  Artinya,
misalkan $(X,d)$ ruang metrik lengkap dan $E \subset X$ himpunan tertutup.
Tunjukkan bahwa $E$ dengan metrik subruang merupakan ruang metrik lengkap.
\end{exercise}

\begin{exercise}
Misalkan $(X,d)$ ruang metrik tak lengkap.  Tunjukkan bahwa terdapat
himpunan tertutup dan terbatas $E \subset X$ yang tidak kompak.
\end{exercise}

\begin{exercise}
Misalkan $(X,d)$ ruang metrik dan $K \subset X$.
Buktikan bahwa $K$ kompak sebagai subhimpunan dari $(X,d)$ jika dan hanya jika $K$
kompak sebagai subhimpunan dirinya sendiri dengan metrik subruang.
\end{exercise}

\begin{exercise} \label{exercise:cauchydepndsonmetric}
Tinjau dua metrik pada $\R$.
Misalkan $d(x,y) \coloneqq \sabs{x-y}$ metrik standar,
dan tetapkan
$d'(x,y) \coloneqq
\babs{ \frac{x}{1+\sabs{x}} - \frac{y}{1+\sabs{y}} }$.
\begin{enumerate}[a)]
\item
Tunjukkan bahwa $(\R,d')$ merupakan ruang metrik (jika Anda telah mengerjakan
\exerciseref{exercise:extendedrealsmetric}, perhitungannya sama).
\item
Tunjukkan bahwa topologinya sama, yaitu bahwa suatu himpunan terbuka dalam
$(\R,d)$ jika dan hanya jika terbuka dalam $(\R,d')$.
\item
Tunjukkan bahwa suatu himpunan kompak dalam
$(\R,d)$ jika dan hanya jika kompak dalam $(\R,d')$.
\item
Tunjukkan bahwa suatu barisan konvergen dalam $(\R,d)$ jika dan hanya jika
barisan tersebut konvergen dalam $(\R,d')$.
\item
Temukan barisan bilangan real yang merupakan barisan Cauchy
dalam $(\R,d')$ tetapi bukan barisan Cauchy dalam $(\R,d)$.
\item
Meskipun $(\R,d)$ lengkap, tunjukkan bahwa $(\R,d')$ tidak lengkap.
\end{enumerate}
\end{exercise}

\begin{exercise} \label{exercise:relativelycompactseq}
\pagebreak[2]
Misalkan $(X,d)$ ruang metrik lengkap.
Himpunan $S \subset X$ disebut \emph{\myindex{relatif kompak}}
jika penutupannya $\widebar{S}$ kompak.
Buktikan bahwa $S \subset X$ relatif kompak jika dan hanya jika
untuk setiap barisan $\{ x_n \}_{n=1}^\infty$ dalam $S$, terdapat subbarisan
$\{ x_{n_k} \}_{k=1}^\infty$ yang konvergen (dalam $X$).
\end{exercise}

%%%%%%%%%%%%%%%%%%%%%%%%%%%%%%%%%%%%%%%%%%%%%%%%%%%%%%%%%%%%%%%%%%%%%%%%%%%%%%

\sectionnewpage
\section{Fungsi kontinu}
\label{sec:metcont}

%mbxINTROSUBSECTION

\sectionnotes{1,5--2 kuliah}

\subsection{Kekontinuan}

\begin{defn}
Misalkan $(X,d_X)$ dan $(Y,d_Y)$ ruang metrik dan $c \in X$.
Fungsi $f \colon X \to Y$ disebut
\emph{kontinu di $c$}\index{kontinu di $c$}
jika untuk setiap $\epsilon > 0$
terdapat $\delta > 0$ sedemikian sehingga apabila $x \in X$ dan $d_X(x,c) <
\delta$, maka
$d_Y\bigl(f(x),f(c)\bigr) < \epsilon$.

\medskip

Jika $f \colon X \to Y$ kontinu di setiap $c \in X$, kita cukup mengatakan
bahwa $f$ adalah \emph{fungsi kontinu}\index{fungsi kontinu!dalam ruang
metrik}\index{fungsi!kontinu}.
\end{defn}

Definisi ini sesuai dengan definisi dalam \chapterref{lim:chapter} ketika
$f$ merupakan fungsi bernilai real pada garis real---tentu saja selama kita
menggunakan metrik standar pada $\R$.

\begin{prop} \label{prop:contiscont}
Misalkan $(X,d_X)$ dan $(Y,d_Y)$ ruang metrik.
Fungsi $f \colon X \to Y$
kontinu di $c \in X$
jika dan hanya jika untuk setiap barisan $\{ x_n \}_{n=1}^\infty$ dalam $X$
yang konvergen ke $c$, barisan $\bigl\{ f(x_n) \bigr\}_{n=1}^\infty$ konvergen
ke $f(c)$.
\end{prop}

\begin{proof}
Andaikan $f$ kontinu di $c$.  Misalkan $\{ x_n \}_{n=1}^\infty$ suatu
barisan dalam $X$ yang konvergen ke $c$.  Untuk setiap $\epsilon > 0$,
terdapat $\delta > 0$ sedemikian sehingga $d_X(x,c) < \delta$ mengakibatkan
$d_Y\bigl(f(x),f(c)\bigr) < \epsilon$.  Jadi, ambil $M$ sedemikian sehingga
untuk semua $n \geq M$, berlaku $d_X(x_n,c) < \delta$, sehingga
$d_Y\bigl(f(x_n),f(c)\bigr) < \epsilon$.  Oleh karena itu, $\bigl\{ f(x_n) \bigr\}_{n=1}^\infty$
konvergen ke $f(c)$.

Sebaliknya, andaikan $f$ tidak kontinu di $c$.
Maka terdapat $\epsilon > 0$
sedemikian sehingga untuk setiap $n \in \N$ terdapat $x_n \in X$
dengan
$d_X(x_n,c) < \nicefrac{1}{n}$ dan $d_Y\bigl(f(x_n),f(c)\bigr) \geq
\epsilon$.  Barisan $\{ x_n \}_{n=1}^\infty$ konvergen ke $c$, tetapi
$\bigl\{ f(x_n) \bigr\}_{n=1}^\infty$
tidak konvergen ke $f(c)$.
\end{proof}

\begin{example}
Misalkan $f \colon \R^2 \to \R$ suatu polinom.  Artinya,
\begin{equation*}
f(x,y) =
\sum_{j=0}^d
\sum_{k=0}^{d-j}
a_{jk}\,x^jy^k =
a_{0\,0} + a_{1\,0} \, x +
a_{0\,1} \, y+  
a_{2\,0} \, x^2+  
a_{1\,1} \, xy+  
a_{0\,2} \, y^2+ \cdots +
a_{0\,d} \, y^d ,
\end{equation*}
untuk suatu $d \in \N$ (derajatnya) dan $a_{jk} \in \R$.  Kita akan menunjukkan 
bahwa $f$ kontinu.  Misalkan $\bigl\{ (x_n,y_n) \bigr\}_{n=1}^\infty$ suatu barisan
dalam $\R^2$ yang konvergen ke $(x,y) \in \R^2$.  Kita telah membuktikan bahwa ini
berarti $\lim_{n\to\infty} x_n = x$ dan $\lim_{n\to\infty} y_n = y$.
Menurut \propref{prop:contalg},
\begin{equation*}
\lim_{n\to\infty}
f(x_n,y_n) =
\lim_{n\to\infty}
\sum_{j=0}^d
\sum_{k=0}^{d-j}
a_{jk} \, x_n^jy_n^k 
=
\sum_{j=0}^d
\sum_{k=0}^{d-j}
a_{jk} \, x^jy^k
=
f(x,y) .
\end{equation*}
Jadi, $f$ kontinu di $(x,y)$, dan karena $(x,y)$ dipilih sembarang, $f$
kontinu di setiap titik.  Dengan cara yang sama,
polinom dalam $n$ variabel juga kontinu.
\end{example}

Berhati-hatilah ketika mengambil limit secara terpisah.
Tinjau $f \colon \R^2 \to \R$ yang didefinisikan oleh $f(x,y) \coloneqq \frac{xy}{x^2+y^2}$
di luar titik asal dan $f(0,0) \coloneqq 0$.
Lihat \figureref{fig:xyxsqysq}.
Dalam \exerciseref{exercise:dicontR2},
Anda diminta membuktikan bahwa $f$ tidak kontinu di titik asal.
Namun, untuk setiap $y$, fungsi $g(x) \coloneqq f(x,y)$
kontinu,
dan untuk setiap $x$, fungsi $h(y) \coloneqq f(x,y)$ kontinu.
\begin{myfigureht}
\myincludegraphics{xyxsqysq}{%
Grafik permukaan suatu fungsi yang tampak terjepit di
titik asal, sehingga jika kita mendekati titik asal pada bidang xy
dari arah yang berbeda-beda, kita dapat mendekati nilai yang berbeda-beda.}
\caption{Grafik $\frac{xy}{x^2+y^2}$.\label{fig:xyxsqysq}}
\end{myfigureht}

\begin{example}
Misalkan $X$ ruang metrik dan $f \colon X \to \C$ fungsi bernilai
kompleks.  Tuliskan $f(p) = g(p) + i \, h(p)$, dengan $g \colon X \to \R$
dan $h \colon X \to \R$ masing-masing bagian real dan imajiner dari $f$.
Maka $f$ kontinu di $c \in X$ jika dan hanya jika bagian real
dan bagian imajinernya kontinu di~$c$.  
Fakta ini berlaku karena $\bigl\{ f(p_n) = g(p_n) + i \, h(p_n) \bigr\}_{n=1}^\infty$
konvergen ke $f(p) = g(p) + i \, h(p)$ jika dan hanya jika
$\bigl\{ g(p_n) \bigr\}_{n=1}^\infty$ konvergen ke $g(p)$ dan
$\bigl\{ h(p_n) \bigr\}_{n=1}^\infty$ konvergen ke $h(p)$.
\end{example}

\subsection{Kekompakan dan kekontinuan}

Pemetaan kontinu tidak selalu memetakan himpunan tertutup ke himpunan tertutup.  Sebagai contoh,
$f \colon (0,1) \to \R$ yang didefinisikan oleh $f(x) \coloneqq x$ memetakan himpunan $(0,1)$, yang
tertutup dalam $(0,1)$, ke himpunan $(0,1)$, yang tidak tertutup dalam $\R$.
Akan tetapi, citra setiap himpunan kompak di bawah pemetaan kontinu tetap kompak.

\begin{lemma} \label{lemma:continuouscompact}
Misalkan $(X,d_X)$ dan $(Y,d_Y)$ ruang metrik
dan $f \colon X \to Y$ suatu fungsi kontinu.  Jika
$K \subset X$ himpunan kompak, maka $f(K)$ himpunan kompak.
\end{lemma}

\begin{proof}
Suatu barisan dalam $f(K)$ dapat dituliskan sebagai
$\bigl\{ f(x_n) \bigr\}_{n=1}^\infty$, dengan
$\{ x_n \}_{n=1}^\infty$ suatu barisan dalam $K$.  Himpunan $K$ kompak,
sehingga terdapat subbarisan
$\{ x_{n_j} \}_{j=1}^\infty$ yang konvergen ke suatu $x \in K$.
Karena fungsi tersebut kontinu,
\begin{equation*}
\lim_{j\to\infty} f(x_{n_j}) = f(x) \in f(K) .
\end{equation*}
Jadi, setiap barisan dalam $f(K)$ mempunyai subbarisan yang konvergen ke
suatu titik dalam $f(K)$, sehingga $f(K)$ kompak menurut \thmref{thm:mscompactisseqcpt}.
\end{proof}

Seperti sebelumnya, $f \colon X \to \R$ mencapai
\emph{\myindex{minimum mutlak}} di $c \in X$ jika
\begin{equation*}
f(x) \geq f(c) \qquad \text{untuk semua } x \in X.
\end{equation*}
Di sisi lain, $f$ mencapai
\emph{\myindex{maksimum mutlak}} di $c \in X$ jika
\begin{equation*}
f(x) \leq f(c) \qquad \text{untuk semua } x \in X.
\end{equation*}

\begin{thm}
Misalkan $(X,d)$ ruang metrik kompak tak kosong
dan misalkan $f \colon X \to \R$ kontinu.  Maka
$f$ terbatas dan, bahkan,
$f$ mencapai minimum mutlak dan maksimum mutlak pada $X$.
\end{thm}

\begin{proof}
Karena $X$ kompak dan $f$ kontinu,
$f(X) \subset \R$ kompak.  Oleh karena itu, $f(X)$ tertutup
dan terbatas.  Khususnya,
$\sup f(X) \in f(X)$ dan
$\inf f(X) \in f(X)$, karena baik supremum maupun infimum
dapat dihampiri oleh barisan dalam $f(X)$ dan $f(X)$ tertutup.
Oleh karena itu, terdapat $x \in X$ sedemikian sehingga $f(x) = \sup f(X)$
dan $y \in X$ sedemikian sehingga $f(y) = \inf f(X)$.
\end{proof}

\subsection{Kekontinuan dan topologi}

Mari kita melihat cara mendefinisikan kekontinuan dalam istilah topologi, yaitu
himpunan-himpunan terbuka.  Kita telah melihat bahwa topologi menentukan barisan mana yang 
konvergen, sehingga tidak mengherankan bahwa topologi juga
menentukan kekontinuan fungsi.

\begin{lemma} \label{lemma:mstopocontloc}
Misalkan $(X,d_X)$ dan $(Y,d_Y)$ ruang metrik.
Fungsi $f \colon X \to Y$ kontinu di $c \in X$
jika dan hanya jika untuk setiap lingkungan terbuka $U$ bagi $f(c)$ dalam $Y$, himpunan
$f^{-1}(U)$ memuat suatu lingkungan terbuka bagi $c$ dalam $X$.
Lihat \figureref{fig:mscontfuncpt}.
\end{lemma}

\begin{myfigureht}
\myincludepdft{mscontfuncpt}{%
Di sisi kiri terdapat daerah berarsir gelap bertanda W dengan batas bertitik-titik,
yang termuat dalam daerah berarsir terang bertanda invers f dari U
dengan batas garis-titik.  Di sisi kanan terdapat daerah berarsir terang
bertanda U dengan batas bertitik-titik.  Sebuah titik bertanda c di daerah W
dihubungkan oleh panah bertanda f ke titik f dari c di daerah U.}
\caption{Untuk setiap lingkungan $U$ bagi $f(c)$, himpunan $f^{-1}(U)$ memuat suatu
lingkungan terbuka $W$ bagi $c$.\label{fig:mscontfuncpt}}
\end{myfigureht}

\begin{proof}
Pertama, andaikan $f$ kontinu di $c$.
Misalkan $U$ suatu lingkungan terbuka bagi $f(c)$
dalam $Y$; maka $B_Y\bigl(f(c),\epsilon\bigr) \subset U$ untuk suatu $\epsilon >
0$.  Oleh kekontinuan $f$, terdapat $\delta > 0$
sedemikian sehingga apabila $x$ memenuhi $d_X(x,c) < \delta$, maka
$d_Y\bigl(f(x),f(c)\bigr) < \epsilon$.  Dengan kata lain,
\begin{equation*}
B_X(c,\delta) \subset f^{-1}\bigl(B_Y\bigl(f(c),\epsilon\bigr)\bigr) \subset
f^{-1}(U) ,
\end{equation*}
dan $B_X(c,\delta)$ merupakan suatu lingkungan terbuka bagi $c$.

Sebaliknya,
misalkan $\epsilon > 0$ diberikan.  Jika
$f^{-1}\bigl(B_Y\bigl(f(c),\epsilon\bigr)\bigr)$ memuat suatu
lingkungan terbuka $W$ bagi $c$, himpunan itu memuat suatu bola.  Artinya, terdapat $\delta > 0$
sedemikian sehingga
\begin{equation*}
B_X(c,\delta) \subset W \subset f^{-1}\bigl(B_Y\bigl(f(c),\epsilon\bigr)\bigr) .
\end{equation*}
Ini tepat berarti bahwa jika $d_X(x,c) < \delta$,
maka $d_Y\bigl(f(x),f(c)\bigr) < \epsilon$.
Jadi, $f$ kontinu di~$c$.
\end{proof}

\begin{thm} \label{thm:mstopocont}
Misalkan $(X,d_X)$ dan $(Y,d_Y)$ ruang metrik.  Fungsi $f \colon X \to Y$
kontinu jika dan hanya jika
untuk setiap himpunan terbuka $U \subset Y$, $f^{-1}(U)$ terbuka dalam $X$.
\end{thm}

Pembuktiannya mengikuti \lemmaref{lemma:mstopocontloc} dan diserahkan sebagai
latihan.

\begin{example}
Misalkan $f \colon X \to Y$ fungsi kontinu.
\thmref{thm:mstopocont} menyatakan bahwa jika $E \subset Y$ tertutup, maka 
$f^{-1}(E) = X \setminus f^{-1}(E^c)$ juga tertutup.  Oleh karena itu, jika
kita mempunyai fungsi kontinu
$f \colon X \to \R$, maka
\emph{\myindex{himpunan nol}} dari $f$, yaitu, 
$f^{-1}(0) = \bigl\{ x \in X : f(x) = 0 \bigr\}$, tertutup.
Kita baru saja membuktikan hasil paling dasar dalam
\emph{\myindex{geometri aljabar}}, yaitu kajian tentang
himpunan nol polinom: Himpunan nol suatu polinom tertutup.

Demikian pula, himpunan tempat $f$ bernilai taknegatif,
$f^{-1}\bigl( [0,\infty) \bigr) = \bigl\{ x \in X : f(x) \geq 0 \bigr\}$,
tertutup.  Di sisi lain,
himpunan tempat $f$ bernilai positif,
$f^{-1}\bigl( (0,\infty) \bigr) = \bigl\{ x \in X : f(x) > 0 \bigr\}$,
terbuka.  
\end{example}

\subsection{Kekontinuan seragam}

Seperti halnya untuk fungsi kontinu
pada garis real, dalam definisi kekontinuan
kadang-kadang berguna untuk dapat memilih
satu $\delta$ bagi semua titik.

\begin{defn}
Misalkan $(X,d_X)$ dan $(Y,d_Y)$ ruang metrik.
Fungsi $f \colon X \to Y$ disebut
\emph{kontinu seragam}\index{kontinu seragam!dalam ruang metrik}
jika untuk setiap $\epsilon > 0$
terdapat $\delta > 0$ sedemikian sehingga untuk setiap $p,q \in X$ dengan $d_X(p,q) <
\delta$, berlaku
$d_Y\bigl(f(p),f(q)\bigr) < \epsilon$.
\end{defn}

Fungsi kontinu seragam pasti kontinu, tetapi kebalikannya tidak selalu
berlaku, seperti telah kita lihat.

\begin{thm} \label{thm:Xcompactfunifcont}
Misalkan $(X,d_X)$ dan $(Y,d_Y)$ ruang metrik.
Andaikan $f \colon X \to Y$ kontinu dan $X$ kompak.  Maka
$f$ kontinu seragam.
\end{thm}

\begin{proof}
Ambil $\epsilon > 0$.  Untuk setiap $c \in X$, pilih $\delta_c > 0$ sedemikian sehingga
$d_Y\bigl(f(x),f(c)\bigr) < \nicefrac{\epsilon}{2}$
untuk setiap
$x \in B(c,\delta_c)$.
%$d_X(x,c) < \delta_c$.
Bola-bola
$B(c,\delta_c)$ menyelimuti $X$, dan ruang $X$ kompak.  
Terapkan \hyperref[ms:lebesgue]{lema selimut Lebesgue} untuk memperoleh 
$\delta > 0$ sedemikian sehingga untuk setiap $x \in X$, terdapat $c \in X$
dengan $B(x,\delta) \subset B(c,\delta_c)$.

Andaikan $p, q \in X$ dengan $d_X(p,q) < \delta$.
Pilih $c \in X$ sedemikian sehingga $B(p,\delta) \subset B(c,\delta_c)$.
Maka $q \in B(c,\delta_c)$.  Berdasarkan ketaksamaan segitiga
dan definisi $\delta_c$,
\begin{equation*}
d_Y\bigl(f(p),f(q)\bigr)
\leq
d_Y\bigl(f(p),f(c)\bigr)
+
d_Y\bigl(f(c),f(q)\bigr)
<
\nicefrac{\epsilon}{2}+
\nicefrac{\epsilon}{2} = \epsilon .  \qedhere
\end{equation*}
\end{proof}

Sebagai penerapan kekontinuan seragam,
kita buktikan suatu kriteria yang berguna bagi
kekontinuan fungsi yang didefinisikan oleh integral.
Misalkan $f(x,y)$ fungsi dua variabel dan definisikan
\begin{equation*}
g(y) \coloneqq \int_a^b f(x,y) \,dx .
\end{equation*}
Pertanyaannya adalah apakah $g$ kontinu.
Sesungguhnya, kita menanyakan kapan dua operasi limit dapat dipertukarkan.
Hal ini tidak selalu dapat dilakukan, sehingga diperlukan suatu hipotesis
tambahan.  Syarat cukup yang berguna (tetapi tidak
perlu) adalah bahwa $f$ kontinu pada persegi panjang tertutup.

\begin{prop} \label{prop:integralcontcont}
Jika $f \colon [a,b] \times [c,d] \to \R$ kontinu,
maka $g \colon [c,d] \to \R$ yang didefinisikan oleh
\begin{equation*}
g(y) \coloneqq \int_a^b f(x,y) \,dx  \qquad \text{bersifat kontinu}.
\end{equation*}
\end{prop}

\begin{proof}
Tetapkan $y \in [c,d]$ dan
ambil $\epsilon > 0$.
Karena $f$ kontinu pada himpunan kompak $[a,b] \times [c,d]$, maka $f$
kontinu seragam.  
Secara khusus, terdapat $\delta > 0$ sedemikian sehingga
untuk setiap $z \in [c,d]$ dengan
$\sabs{z-y} < \delta$, berlaku
$\babs{f(x,z)-f(x,y)} < \frac{\epsilon}{b-a}$ untuk semua $x \in [a,b]$.
Jadi, andaikan $\sabs{z-y} < \delta$.  Maka
\begin{multline*}
\babs{
g(z)-
g(y)
}
=
\abs{
\int_a^b 
f(x,z) \,dx 
-
\int_a^b 
f(x,y) \,dx 
}
\\
=
\abs{
\int_a^b 
\bigl(
f(x,z) - f(x,y)
\bigr)
\,dx 
}
\leq
(b-a)
\frac{\epsilon}{b-a}
= \epsilon . \qedhere
\end{multline*}
\end{proof}

Dalam penerapan, jika kita tertarik pada kekontinuan di $y_0$, kita hanya
perlu menerapkan proposisi tersebut pada $[a,b] \times [y_0-\epsilon,y_0+\epsilon]$
untuk suatu $\epsilon > 0$ yang kecil.  Sebagai contoh, jika $f$ kontinu pada
$[a,b] \times \R$, maka $g$ kontinu pada $\R$.

\begin{example}
Fungsi
\emph{kontinu Lipschitz}%
\index{kontinu Lipschitz!dalam ruang metrik}%
\index{fungsi!Lipschitz}
merupakan contoh lain fungsi kontinu seragam yang berguna.  Artinya, jika
$(X,d_X)$ dan $(Y,d_Y)$ ruang metrik, maka $f \colon X \to Y$
disebut Lipschitz atau $K$-Lipschitz jika terdapat $K \in \R$ sedemikian sehingga
\begin{equation*}
d_Y\bigl(f(p),f(q)\bigr) \leq K \, d_X(p,q)
\qquad \text{untuk semua } p,q \in X.
\end{equation*}
Fungsi Lipschitz bersifat kontinu seragam:
Ambil $\delta = \nicefrac{\epsilon}{K}$ (dengan mengandaikan $K > 0$).
Suatu fungsi dapat kontinu seragam
tetapi tidak Lipschitz,
seperti telah kita lihat: $\sqrt{x}$ pada $[0,1]$
kontinu seragam tetapi tidak Lipschitz.

Perlu disebutkan bahwa
jika suatu fungsi bersifat Lipschitz, biasanya paling mudah
untuk langsung menunjukkan bahwa fungsi tersebut Lipschitz meskipun kita hanya
ingin mengetahui kekontinuannya.
\end{example}

\subsection{Titik akumulasi dan limit fungsi}

Meskipun kita tidak memulai pembahasan kekontinuan dengan titik akumulasi dan
belum akan membutuhkannya sampai
\volIIref{Jilid II\@,}{jauh kemudian,}
mari kita juga
terjemahkan gagasan limit fungsi dari garis real ke ruang
metrik.
Sekali lagi, kita perlu memulai dengan titik akumulasi.

\begin{defn}
Misalkan $(X,d)$ ruang metrik dan
$S \subset X$. Titik $p \in X$ disebut
\emph{titik akumulasi}\index{titik akumulasi!dalam ruang metrik} dari $S$
jika untuk setiap $\epsilon > 0$, himpunan $B(p,\epsilon) \cap S
\setminus \{ p \}$ tidak kosong.
\end{defn}

Tidak cukup bahwa $p$ berada dalam penutupan $S$;
titik tersebut harus berada dalam penutupan
$S \setminus \{ p \}$ (latihan).
Jadi, $p$ merupakan titik akumulasi jika dan hanya jika terdapat barisan dalam $S \setminus
\{ p \}$ yang konvergen ke $p$.

\begin{defn}
\index{limit!fungsi dalam ruang metrik}%
Misalkan $(X,d_X)$ dan $(Y,d_Y)$ ruang metrik, $S \subset X$, $p \in X$ titik akumulasi dari $S$,
dan $f \colon S \to Y$ suatu fungsi.
Andaikan terdapat $L \in Y$ dan, untuk setiap $\epsilon > 0$,
terdapat $\delta > 0$ sedemikian sehingga untuk setiap $x \in S \setminus \{ p \}$
dengan $d_X(x,p) < \delta$, berlaku
\begin{equation*}
d_Y\bigl(f(x),L\bigr) < \epsilon .
\end{equation*}
Kita katakan bahwa $f(x)$
\emph{konvergen}\index{konvergen!fungsi dalam ruang metrik} ke $L$ saat $x$ menuju
$p$, dan $L$ merupakan \emph{limit} dari $f(x)$ saat $x$
menuju $p$.  Jika $L$ tunggal, kita tulis
\glsadd{not:limfunc}
\begin{equation*}
\lim_{x \to p} f(x) \coloneqq L .
\end{equation*}
Jika $f(x)$ tidak konvergen saat $x$ menuju $p$, kita katakan bahwa $f$
\emph{divergen}\index{divergen!fungsi dalam ruang metrik} di $p$.
\end{defn}

Seperti biasa, kita buktikan bahwa limit tersebut tunggal jika ada.
Pembuktiannya merupakan terjemahan langsung dari pembuktian
dalam \chapterref{lim:chapter}, sehingga kita serahkan sebagai latihan.

\begin{prop} \label{prop:mslimitisunique}
Misalkan $(X,d_X)$ dan $(Y,d_Y)$ ruang metrik, $S \subset X$, $p \in X$
titik akumulasi dari $S$, dan $f \colon S \to Y$ suatu fungsi
sedemikian sehingga $f(x)$ konvergen saat $x$ menuju $p$.  Maka
limit dari $f(x)$ saat $x$ menuju $p$ bersifat tunggal.
\end{prop}

Dalam setiap ruang metrik, seperti dalam $\R$, limit fungsi dapat
diganti dengan limit sekuensial.  Pembuktiannya sekali lagi merupakan terjemahan langsung
dari pembuktian dalam \chapterref{lim:chapter}, dan kita serahkan sebagai
latihan.  Dengan demikian, sesungguhnya kita hanya perlu membuktikan berbagai hal untuk
limit sekuensial.

\begin{lemma}\label{ms:seqflimit:lemma}
Misalkan $(X,d_X)$ dan $(Y,d_Y)$ ruang metrik, $S \subset X$, $p \in X$
titik akumulasi dari $S$, dan $f \colon S \to Y$ suatu fungsi.

Maka
$f(x)$ konvergen ke $L \in Y$ saat $x$ menuju $p$ jika dan hanya jika, untuk setiap
barisan $\{ x_n \}_{n=1}^\infty$
dalam $S \setminus \{p\}$
sedemikian sehingga $\lim_{n\to\infty} x_n = p$,
barisan $\bigl\{ f(x_n) \bigr\}_{n=1}^\infty$ konvergen ke $L$.
\end{lemma}

Dengan menerapkan \propref{prop:contiscont} atau definisinya secara langsung, kita peroleh
(latihan), seperti dalam \chapterref{lim:chapter}, bahwa untuk titik akumulasi $p \in S$ dari $S
\subset X$, fungsi
$f \colon S \to Y$ kontinu di $p$ jika dan hanya jika
\begin{equation*}
\lim_{x \to p} f(x) = f(p) .
\end{equation*}

\subsection{Latihan}

\begin{exercise}
Tinjau $\N \subset \R$ dengan metrik standar.  Misalkan $(X,d)$ suatu
ruang metrik dan $f \colon X \to \N$ fungsi kontinu.
\begin{enumerate}[a)]
\item
Buktikan bahwa jika $X$ terhubung,
maka $f$ konstan (citra $f$ hanya terdiri atas satu nilai).
\item
Temukan contoh ketika $X$ tidak terhubung dan $f$ tidak konstan.
\end{enumerate}
\end{exercise}

\begin{exercise} \label{exercise:dicontR2}
\pagebreak[2]
Definisikan $f \colon \R^2 \to \R$ dengan $f(0,0) \coloneqq 0$, dan
$f(x,y) \coloneqq \frac{xy}{x^2+y^2}$ jika $(x,y) \neq (0,0)$.
Lihat \figureref{fig:xyxsqysq}.
\begin{enumerate}[a)]
\item
Tunjukkan bahwa untuk setiap $x$ tetap,
fungsi yang memetakan $y$ ke $f(x,y)$ kontinu.  Demikian pula,
untuk setiap $y$ tetap, fungsi yang memetakan $x$ ke $f(x,y)$ kontinu.
\item
Tunjukkan bahwa $f$ tidak kontinu.
\end{enumerate}
\end{exercise}

\begin{exercise} 
\pagebreak[2]
Misalkan $(X,d_X)$, $(Y,d_Y)$ ruang metrik dan
$f \colon X \to Y$ kontinu.
Misalkan $A \subset X$.
\begin{enumerate}[a)]
\item
Tunjukkan bahwa $f(\widebar{A}) \subset \overline{f(A)}$.
\item
Tunjukkan melalui contoh bahwa $f(\widebar{A})$ dapat merupakan subhimpunan sejati dari
$\overline{f(A)}$.
\end{enumerate}
\end{exercise}

\begin{exercise}
Buktikan \thmref{thm:mstopocont}.  Petunjuk: Gunakan \lemmaref{lemma:mstopocontloc}.
\end{exercise}

\begin{exercise} \label{exercise:msconnconn}
Misalkan $f \colon X \to Y$ kontinu untuk ruang metrik $(X,d_X)$
dan $(Y,d_Y)$.  Tunjukkan bahwa jika $X$ terhubung, maka $f(X)$ terhubung.
\end{exercise}

\begin{exercise}
Buktikan versi berikut dari
\hyperref[IVT:thm]{teorema nilai antara}.  Misalkan $(X,d)$ ruang metrik
terhubung dan $f \colon X \to \R$ fungsi kontinu.  Andaikan 
$x_0,x_1 \in X$ dan $y \in \R$ sedemikian sehingga $f(x_0) < y < f(x_1)$.
Buktikan bahwa terdapat $z \in X$ sedemikian sehingga $f(z) = y$.
Petunjuk: Lihat \exerciseref{exercise:msconnconn}.
\end{exercise}

\begin{exercise}
Fungsi kontinu $f \colon X \to Y$ antara ruang metrik $(X,d_X)$ dan
$(Y,d_Y)$ disebut \emph{\myindex{proper}}
jika untuk setiap himpunan kompak $K \subset Y$, himpunan $f^{-1}(K)$ kompak.
Andaikan fungsi kontinu $f \colon (0,1) \to (0,1)$ proper dan
$\{ x_n \}_{n=1}^\infty$ barisan dalam $(0,1)$ yang konvergen ke $0$.  Tunjukkan bahwa
$\bigl\{ f(x_n) \bigr\}_{n=1}^\infty$ tidak memiliki subbarisan yang konvergen dalam $(0,1)$.
\end{exercise}

\begin{exercise}
Misalkan $(X,d_X)$ dan $(Y,d_Y)$ ruang metrik dan
$f \colon X \to Y$ fungsi kontinu yang injektif dan surjektif.  Andaikan
$X$ kompak.  Buktikan bahwa invers $f^{-1} \colon Y \to X$
kontinu.
\end{exercise}

\begin{exercise}
Tinjau ruang metrik fungsi kontinu $C\bigl([0,1],\R\bigr)$.  Misalkan
$k \colon [0,1] \times [0,1] \to \R$ fungsi kontinu.
Untuk setiap $f \in C\bigl([0,1],\R\bigr)$, definisikan
\begin{equation*}
\varphi_f(x) \coloneqq \int_0^1 k(x,y) f(y) \,dy .
\end{equation*}
\begin{enumerate}[a)]
\item
Tunjukkan bahwa $T(f) \coloneqq \varphi_f$ mendefinisikan fungsi
$T \colon C\bigl([0,1],\R\bigr) \to C\bigl([0,1],\R\bigr)$.
\item
Tunjukkan bahwa $T$ kontinu.
\end{enumerate}
\end{exercise}

\begin{samepage}
\begin{exercise}
Misalkan $(X,d)$ ruang metrik.
\begin{enumerate}[a)]
\item
Jika $p \in X$,
tunjukkan bahwa $f \colon X \to \R$ yang didefinisikan
oleh $f(x) \coloneqq d(x,p)$ kontinu.
\item
Definisikan metrik pada $X \times X$ seperti dalam \exerciseref{exercise:mscross} bagian
b, dan tunjukkan bahwa $g \colon X \times X \to \R$ yang didefinisikan oleh
$g(x,y) \coloneqq d(x,y)$ kontinu.
\item
Tunjukkan bahwa jika $K_1$ dan $K_2$ subhimpunan kompak tak kosong dari $X$, maka
terdapat $p \in K_1$ dan $q \in K_2$ sedemikian sehingga $d(p,q)$ bernilai minimum,
yaitu, $d(p,q) = \inf \{ d(x,y) \colon x \in K_1, y \in K_2 \}$.
\end{enumerate}
\end{exercise}
\end{samepage}

\begin{exercise}
Misalkan $(X,d)$ ruang metrik kompak tak kosong, dan misalkan $C(X,\R)$ himpunan
fungsi kontinu bernilai real.  Definisikan
\begin{equation*}
d(f,g) \coloneqq \snorm{f-g}_X = \sup_{x \in X} \babs{f(x)-g(x)} .
\end{equation*}
\begin{enumerate}[a)]
\item
Tunjukkan bahwa $d$ menjadikan $C(X,\R)$ ruang metrik.
\item
Tunjukkan bahwa untuk setiap $x \in X$, fungsi evaluasi
$E_x \colon C(X,\R) \to \R$ yang didefinisikan oleh $E_x(f) \coloneqq f(x)$
merupakan fungsi kontinu.
\end{enumerate}
\end{exercise}

\begin{samepage}
\begin{exercise}
Misalkan $C\bigl([a,b],\R\bigr)$ himpunan fungsi kontinu dan
$C^1\bigl([a,b],\R\bigr)$\glsadd{not:contdifffuncs} himpunan fungsi yang
diferensiabel dan memiliki turunan kontinu pada $[a,b]$.
Definisikan
\begin{equation*}
d_{C}(f,g) \coloneqq \snorm{f-g}_{[a,b]}
\qquad \text{dan} \qquad
d_{C^1}(f,g) \coloneqq \snorm{f-g}_{[a,b]} + \snorm{f'-g'}_{[a,b]},
\end{equation*}
di mana $\snorm{\cdot}_{[a,b]}$ adalah norma seragam.
Berdasarkan \exampleref{example:msC01} dan \exerciseref{exercise:C1ab}, kita mengetahui bahwa
$C\bigl([a,b],\R\bigr)$ dengan $d_C$ merupakan ruang metrik, dan
begitu pula
$C^1\bigl([a,b],\R\bigr)$ dengan $d_{C^1}$.
\begin{enumerate}[a)]
\item
Buktikan bahwa operator diferensiasi $D \colon 
C^1\bigl([a,b],\R\bigr) \to C\bigl([a,b],\R\bigr)$ yang didefinisikan oleh
$D(f) \coloneqq f'$ kontinu.
\item
Sebaliknya, jika kita menggunakan metrik $d_C$ pada $C^1\bigl([a,b],\R\bigr)$,
buktikan bahwa operator diferensiasi tidak lagi kontinu.  Petunjuk: Tinjau
$\sin(n x)$.
\end{enumerate}
\end{exercise}
\end{samepage}

\begin{exercise}
Misalkan $(X,d)$ ruang metrik, $S \subset X$, dan $p \in X$.  Buktikan bahwa
$p$ merupakan titik akumulasi dari $S$ jika dan hanya jika $p \in \overline{S \setminus \{
p \}}$.
\end{exercise}

\begin{exercise}
Buktikan \propref{prop:mslimitisunique}.
\end{exercise}

\begin{exercise}
Buktikan \lemmaref{ms:seqflimit:lemma}.
\end{exercise}

\begin{exercise}
Misalkan $(X,d_X)$ dan $(Y,d_Y)$ ruang metrik, $S \subset X$, $p \in S$
suatu titik akumulasi dari $S$, dan misalkan $f \colon S \to Y$ suatu fungsi.
Buktikan bahwa
$f \colon S \to Y$ kontinu di $p$ jika dan hanya jika
\begin{equation*}
\lim_{x \to p} f(x) = f(p) .
\end{equation*}
\end{exercise}

\begin{exercise}
Definisikan
\begin{equation*}
f(x,y) \coloneqq
\begin{cases}
\frac{2xy}{x^4+y^2} & \text{jika } (x,y) \neq (0,0), \\
0 & \text{jika } (x,y) = (0,0) .
\end{cases}
\end{equation*}
\begin{enumerate}[a)]
\item
Tunjukkan bahwa untuk setiap $y$ tetap, fungsi yang memetakan $x$ ke $f(x,y)$
kontinu dan karenanya terintegralkan Riemann.
\item
Tunjukkan bahwa untuk setiap $x$ tetap, fungsi yang memetakan $y$ ke $f(x,y)$ kontinu.
\item
Tunjukkan bahwa $f$ tidak kontinu di $(0,0)$.
\item
Sekarang tunjukkan bahwa $g(y) \coloneqq \int_0^1 f(x,y)\,dx$ tidak kontinu di $y=0$.
\end{enumerate}
Catatan: Anda boleh menggunakan pengetahuan tentang $\arctan$ dari kalkulus,
khususnya bahwa $\frac{d}{ds} \bigl[ \arctan(s) \bigr] = \frac{1}{1+s^2}$.
\end{exercise}

\begin{exercise} \label{exercise:integralcontcontextra}
Buktikan versi yang lebih kuat dari \propref{prop:integralcontcont}:
Jika $f \colon (a,b) \times (c,d) \to \R$ fungsi kontinu yang terbatas,
maka $g \colon (c,d) \to \R$ yang didefinisikan oleh
\begin{equation*}
g(y) \coloneqq \int_a^b f(x,y) \,dx  \qquad \text{bersifat kontinu}.
\end{equation*}
Petunjuk: Pertama-tama, integralkan pada $[a+\nicefrac{1}{n},b-\nicefrac{1}{n}]$.
\end{exercise}

%%%%%%%%%%%%%%%%%%%%%%%%%%%%%%%%%%%%%%%%%%%%%%%%%%%%%%%%%%%%%%%%%%%%%%%%%%%%%%

\sectionnewpage
\section{Teorema titik tetap dan teorema Picard sekali lagi}
\label{sec:metpicard}

%mbxINTROSUBSECTION

\sectionnotes{1 kuliah (opsional, tidak memerlukan \sectionref{sec:picard})}

Dalam bagian ini kita buktikan teorema titik tetap untuk pemetaan
kontraksi.  Sebagai penerapan, kita buktikan teorema Picard, yang telah kita buktikan
tanpa ruang metrik dalam \sectionref{sec:picard}.
Pembuktian yang disajikan di sini serupa, tetapi berlangsung jauh
lebih lancar dengan ruang metrik dan teorema titik tetap.

\subsection{Teorema titik tetap}

\begin{defn}
Misalkan $(X,d_X)$ dan $(Y,d_Y)$ ruang metrik.
Pemetaan
$\varphi \colon X \to Y$ disebut \emph{\myindex{kontraksi}}
(atau pemetaan kontraktif) jika merupakan
pemetaan $k$-Lipschitz untuk suatu $0 \leq k < 1$, yaitu jika terdapat $0 \leq k < 1$ sedemikian sehingga
\begin{equation*}
d_Y\bigl(\varphi(p),\varphi(q)\bigr) \leq k\, d_X(p,q)
\qquad \text{untuk semua } p,q \in X.
\end{equation*}

\medskip

Diberikan pemetaan $\varphi \colon X \to X$, titik $x \in X$ disebut
\emph{\myindex{titik tetap}}
jika $\varphi(x)=x$.
\end{defn}

\begin{thm}%
[Prinsip pemetaan kontraksi\index{prinsip pemetaan kontraksi}
atau \myindex{teorema titik tetap Banach}\index{teorema titik tetap}%
\footnote{Dinamai menurut matematikawan Polandia
\href{https://en.wikipedia.org/wiki/Stefan_Banach}{Stefan Banach}
(1892--1945), yang pertama kali menyatakan teorema ini pada tahun 1922.}]
\label{thm:contr}
Misalkan $(X,d)$ ruang metrik lengkap tak kosong dan $\varphi \colon X \to X$ suatu
kontraksi.
Maka $\varphi$ memiliki tepat satu titik tetap.
\end{thm}

Syarat \emph{lengkap} dan \emph{kontraksi} diperlukan.
Lihat \exerciseref{exercise:nofixedpoint}.

\begin{proof}
Pilih $x_0 \in X$.
Definisikan barisan $\{ x_n \}_{n=0}^\infty$ dengan $x_{n+1} \coloneqq \varphi(x_n)$.  Maka
\begin{equation*}
d(x_{n+1},x_n) = d\bigl(\varphi(x_n),\varphi(x_{n-1})\bigr)
\leq k d(x_n,x_{n-1}) .
%\leq \cdots
%\leq k^n d(x_1,x_0) .
\end{equation*}
Dengan mengulang langkah ini $n$ kali, diperoleh $d(x_{n+1},x_n) \leq k^n d(x_1,x_0)$.
Untuk $m > n$,
\begin{equation*}
\begin{split}
d(x_m,x_n)
& \leq \sum_{i=n}^{m-1} d(x_{i+1},x_i) \\
& \leq \sum_{i=n}^{m-1} k^i d(x_1,x_0) \\
& = k^n d(x_1,x_0) \sum_{i=0}^{m-n-1} k^i \\
& \leq k^n d(x_1,x_0) \sum_{i=0}^{\infty} k^i
= k^n d(x_1,x_0) \frac{1}{1-k} .
\end{split}
\end{equation*}
Secara khusus, barisan tersebut merupakan barisan Cauchy (mengapa?).  Karena $X$ lengkap,
tetapkan $x \coloneqq \lim_{n\to\infty} x_n$, dan kita nyatakan bahwa $x$
merupakan titik tetap tunggal yang dicari.

Titik tetap?  Fungsi $\varphi$ merupakan kontraksi,
sehingga kontinu Lipschitz:
\begin{equation*}
\varphi(x) = \varphi\Bigl( \lim_{n\to\infty} x_n\Bigr) = \lim_{n\to\infty}
\varphi(x_n) =
\lim_{n\to\infty} x_{n+1} = x .
\end{equation*}
Tunggal?  Misalkan $x$ dan $y$ titik tetap.
\begin{equation*}
d(x,y) = d\bigl(\varphi(x),\varphi(y)\bigr) \leq k\, d(x,y) .
\end{equation*}
Karena $k < 1$, ketaksamaan tersebut mengakibatkan $d(x,y) = 0$, sehingga $x=y$.  Teorema ini
terbukti.
\end{proof}

Pembuktian ini bersifat konstruktif.  Kita bukan hanya mengetahui bahwa 
terdapat tepat satu titik tetap, melainkan juga mengetahui cara mencarinya.  Mulailah dari
sebarang titik $x_0 \in X$, lalu hitung berturut-turut $\varphi(x_0)$,
$\varphi(\varphi(x_0))$,
$\varphi(\varphi(\varphi(x_0)))$, dan seterusnya untuk memperoleh hampiran yang semakin baik.
Kita bahkan dapat menentukan seberapa jauh
hampiran tersebut dari titik tetap.
Lihat latihan.  Oleh karena itu, gagasan pembuktian ini
berguna dalam penerapan dunia nyata.

\subsection{Teorema Picard}

Kita mulai dengan ruang metrik tempat kita akan menerapkan
teorema titik tetap:
ruang $C\bigl([a,b],\R\bigr)$ dari \exampleref{example:msC01},
yaitu ruang fungsi kontinu $f \colon [a,b] \to \R$ dengan metrik
\begin{equation*}
d(f,g) \coloneqq \snorm{f-g}_{[a,b]} = \sup_{x \in [a,b]} \babs{f(x)-g(x)} .
\end{equation*}
Konvergensi dalam metrik ini adalah konvergensi dalam norma seragam, atau dengan
kata lain, konvergensi seragam. Oleh karena itu,
$C\bigl([a,b],\R\bigr)$ merupakan ruang metrik lengkap;
lihat \propref{prop:CabRcomplete}.

\medskip

Sekarang perhatikan persamaan diferensial biasa
\begin{equation*}
\frac{dy}{dx} = F(x,y) .
\end{equation*}
Diberikan suatu $x_0, y_0$, kita menginginkan fungsi $y=f(x)$ sedemikian sehingga
$f(x_0) = y_0$ dan
\begin{equation*}
f'(x) = F\bigl(x,f(x)\bigr) .
\end{equation*}
Untuk menghindari keharusan memperkenalkan banyak nama, kita sering cukup menulis $y' = F(x,y)$
untuk persamaannya
dan $y(x)$ untuk solusinya.

Contoh paling sederhana adalah persamaan $y' = y$, $y(0) = 1$.
Solusinya adalah fungsi eksponensial $y(x) = e^x$. Contoh yang agak lebih rumit
adalah $y' = -2xy$, $y(0) = 1$, yang solusinya berupa fungsi Gauss
$y(x) = e^{-x^2}$.

Salah satu persoalan halus adalah seberapa lama solusi itu ada.
Perhatikan persamaan $y' = y^2$, $y(0)=1$. Maka $y(x) = \frac{1}{1-x}$ merupakan
solusi. Walaupun $F$ merupakan fungsi yang cukup \myquote{baik} dan khususnya
terdefinisi untuk setiap $x$ dan $y$, solusi tersebut \myquote{meledak} pada $x=1$.
Untuk contoh lain yang berkaitan dengan teorema Picard, lihat \sectionref{sec:picard}.

Kita akan mencari solusi di 
$C\bigl([a,b],\R\bigr)$, yang mungkin terasa aneh pada awalnya karena
kita mencari fungsi yang terdiferensialkan.
Penjelasannya adalah bahwa kita meninjau
persamaan integral yang bersesuaian,
\begin{equation*}
f(x)
=
y_0 + \int_{x_0}^x F\bigl(t,f(t)\bigr)\,dt .
\end{equation*}
Untuk menyelesaikan persamaan integral ini, kita hanya memerlukan fungsi kontinu, sehingga
dalam arti tertentu tugas kita semestinya lebih mudah---terdapat lebih banyak calon fungsi
yang dapat dicoba. Cara berpikir seperti ini cukup lazim dalam menyelesaikan persamaan
diferensial.

\begin{samepage}
\begin{thm}[Teorema keberadaan dan ketunggalan Picard]%
\index{teorema keberadaan dan ketunggalan}\index{teorema Picard}
Misalkan $I, J \subset \R$ merupakan interval tertutup dan terbatas,
misalkan $I^\circ$ dan $J^\circ$ menyatakan interiornya, dan 
misalkan $(x_0,y_0) \in I^\circ \times J^\circ$.
Andaikan $F \colon I \times J \to \R$ kontinu
dan bersifat Lipschitz terhadap variabel kedua, yaitu terdapat
$L \in \R$ sedemikian sehingga
\begin{equation*}
\babs{F(x,y) - F(x,z)} \leq L \sabs{y-z}
\qquad \text{untuk setiap } y,z \in J, x \in I .
\end{equation*}
Maka terdapat $h > 0$ sedemikian sehingga $[x_0-h,x_0+h] \subset I$ dan terdapat tepat satu fungsi terdiferensialkan
$f \colon [x_0 - h, x_0 + h] \to J \subset \R$ sedemikian sehingga
\begin{equation*}
f'(x) = F\bigl(x,f(x)\bigr) \qquad \text{dan} \qquad f(x_0) = y_0.
\end{equation*}
\end{thm}
\end{samepage}

\begin{proof}
Tanpa mengurangi keumuman, andaikan $x_0 =0$ (latihan).
Karena $I \times J$ kompak dan
$F$ kontinu, maka $F$ terbatas.
Jadi, ambil $M > 0$ sedemikian sehingga
$\babs{F(x,y)} \leq M$ untuk setiap $(x,y) \in I\times J$.
Pilih $\alpha > 0$ sedemikian sehingga
$[-\alpha,\alpha] \subset I$ dan $[y_0-\alpha, y_0 + \alpha] \subset J$.
Definisikan
\begin{equation*}
h \coloneqq \min \left\{ \alpha, \frac{\alpha}{M+L\alpha} \right\} .
\end{equation*}
Perhatikan bahwa $[-h,h] \subset I$. Definisikan
\begin{equation*}
Y \coloneqq \Bigl\{ f \in C\bigl([-h,h],\R\bigr) : f\bigl([-h,h]\bigr) \subset J \Bigr\} .
\end{equation*}
Jadi, $Y$ adalah himpunan fungsi kontinu pada $[-h,h]$ yang nilainya berada di
$J$; dengan kata lain, himpunan ini persis terdiri atas
fungsi-fungsi yang membuat $F\bigl(x,f(x)\bigr)$ terdefinisi.
Sebagai latihan, tunjukkan bahwa $Y$ merupakan himpunan bagian tertutup dari $C\bigl([-h,h],\R\bigr)$
(karena $J$ tertutup).
Ruang $C\bigl([-h,h],\R\bigr)$ lengkap, dan
himpunan bagian tertutup dari ruang metrik lengkap merupakan ruang metrik lengkap dengan
metrik subruang; lihat \propref{prop:closedcomplete}. Jadi, $Y$ dengan
metrik subruang merupakan ruang metrik lengkap.
Kita tulis $d(f,g) = \snorm{f-g}_{[-h,h]}$ untuk metrik ini.

Definisikan pemetaan
$T \colon Y \to C\bigl([-h,h],\R\bigr)$ dengan
\begin{equation*}
T(f)(x)
\coloneqq
y_0 + \int_0^x F\bigl(t,f(t)\bigr)\,dt .
\end{equation*}
Sebagai latihan, periksalah bahwa
$T$ terdefinisi dengan baik, dan bahwa untuk $f \in Y$, $T(f)$ benar-benar berada di $C\bigl([-h,h],\R\bigr)$.
Misalkan $f \in Y$ dan $\sabs{x} \leq h$.
Karena $F$ dibatasi oleh $M$, kita peroleh
\begin{equation*}
\begin{split}
\babs{T(f)(x) - y_0}
&= \abs{\int_0^x F\bigl(t,f(t)\bigr)\,dt} \\
& \leq 
\sabs{x} M \leq hM \leq \frac{\alpha M}{M+ L\alpha} \leq \alpha .
\end{split}
\end{equation*}
Jadi, $T(f)\bigl([-h,h]\bigr) \subset [y_0-\alpha,y_0+\alpha] \subset J$, dan
$T(f) \in Y$. Dengan kata lain, $T(Y) \subset Y$. Mulai sekarang,
kita memandang $T$ sebagai pemetaan dari $Y$ ke $Y$.

Akan kita tunjukkan bahwa $T \colon Y \to Y$ merupakan kontraksi. Pertama, untuk $x \in [-h,h]$
dan $f,g \in Y$, kita mempunyai
\begin{equation*}
\Babs{F\bigl(x,f(x)\bigr) - F\bigl(x,g(x)\bigr)} \leq
L\babs{f(x)- g(x)} \leq L \, d(f,g) .
\end{equation*}
Oleh karena itu,
\begin{equation*}
\begin{split}
\babs{T(f)(x) - T(g)(x)}
&= \abs{\int_0^x \Bigl( F\bigl(t,f(t)\bigr) - F\bigl(t,g(t)\bigr)\Bigr)\,dt} \\
& \leq \sabs{x} L \, d(f,g)
 \leq h L\, d(f,g)
 \leq \frac{L\alpha}{M+L\alpha} \, d(f,g) .
\end{split}
\end{equation*}
Kita memilih $M > 0$, sehingga
$\frac{L\alpha}{M+L\alpha} < 1$.
Dengan mengambil supremum atas $x \in [-h,h]$ pada ruas kiri di atas, diperoleh
$d\bigl(T(f),T(g)\bigr) \leq \frac{L\alpha}{M+L\alpha} \, d(f,g)$,
yaitu $T$ merupakan kontraksi.

Teorema titik tetap (\thmref{thm:contr})
memberikan tepat satu $f \in Y$ sedemikian sehingga $T(f) = f$.
Dengan kata lain,
\begin{equation*} %\label{equation:msinteqpicard}
f(x) = y_0 + \int_0^x F\bigl(t,f(t)\bigr)\,dt .
\end{equation*}
Jelas bahwa $f(0) = y_0$.
Menurut teorema dasar kalkulus (\thmref{thm:FTCv2}),
$f$ terdiferensialkan dan turunannya adalah
$F\bigl(x,f(x)\bigr)$.
Fungsi terdiferensialkan bersifat kontinu, sehingga
$f$ adalah satu-satunya fungsi terdiferensialkan $f \colon [-h,h] \to J$
sedemikian sehingga
 $f'(x) = F\bigl(x,f(x)\bigr)$ dan $f(0) = y_0$.
\end{proof}

\subsection{Latihan}

\begin{exnote}
Untuk latihan lain yang berkaitan dengan teorema Picard, lihat \sectionref{sec:picard}.
\end{exnote}

\begin{exercise}
Misalkan $J$ interval tertutup dan terbatas, serta
$Y \coloneqq \bigl\{ f \in C\bigl([-h,h],\R\bigr) : f\bigl([-h,h]\bigr) \subset J \bigr\}$.
Buktikan bahwa $Y \subset C\bigl([-h,h],\R\bigr)$ tertutup.  Petunjuk: $J$ tertutup.
\end{exercise}

\begin{exercise}
Dalam pembuktian teorema Picard,
buktikan bahwa jika $f \colon [-h,h] \to J$ kontinu, maka $F\bigl(t,f(t)\bigr)$
kontinu pada $[-h,h]$ sebagai fungsi terhadap $t$.  Gunakan fakta ini untuk menunjukkan bahwa
\begin{equation*}
T(f)(x)
\coloneqq
y_0 + \int_0^x F\bigl(t,f(t)\bigr)\,dt
\end{equation*}
terdefinisi dengan baik dan bahwa $T(f) \in C\bigl([-h,h],\R\bigr)$.
\end{exercise}


\begin{exercise}
Buktikan bahwa dalam pembuktian teorema Picard,
pernyataan \myquote{Tanpa mengurangi keumuman, andaikan $x_0 = 0$}
dapat dibenarkan.  Artinya, buktikan bahwa jika kita mengetahui teorema tersebut untuk $x_0 = 0$,
teorema itu berlaku sebagaimana dinyatakan.
\end{exercise}

\begin{exercise}
\pagebreak[2]
Misalkan $F \colon \R \to \R$ didefinisikan oleh
$F(x) \coloneqq kx + b$ dengan $0 < k < 1$, $b \in \R$.
\begin{enumerate}[a)]
\item
Buktikan bahwa $F$ merupakan kontraksi.
\item
Tentukan titik tetapnya dan tunjukkan secara langsung bahwa titik itu tunggal.
\end{enumerate}
\end{exercise}

\begin{exercise}
\pagebreak[2]
Misalkan $f \colon [0,\nicefrac{1}{4}] \to [0,\nicefrac{1}{4}]$ didefinisikan oleh
$f(x) \coloneqq x^2$.
\begin{enumerate}[a)]
\item
Buktikan bahwa $f$
merupakan kontraksi, dan tentukan konstanta $k$ terbaik (terkecil) yang memenuhi definisi tersebut.
\item
Tentukan titik tetapnya dan tunjukkan secara langsung bahwa titik itu tunggal.
\end{enumerate}
\end{exercise}

\begin{samepage}
\begin{exercise} \label{exercise:nofixedpoint}
\leavevmode
\begin{enumerate}[a)]
\item
Berikan contoh kontraksi $f \colon X \to X$
pada ruang metrik tak lengkap $X$ yang tidak memiliki
titik tetap.
\item
Berikan pemetaan 1-Lipschitz $f \colon X \to X$ pada ruang metrik lengkap $X$ yang tidak memiliki titik tetap.
\end{enumerate}
\end{exercise}
\end{samepage}

\begin{exercise}
Tinjau $y' =y^2$, $y(0)=1$.  Gunakan skema iterasi
dari pembuktian prinsip pemetaan kontraksi.
Mulailah dengan $f_0(x) = 1$.  Tentukan beberapa 
iterat (sekurang-kurangnya sampai $f_2$).  Buktikan bahwa
limit titik demi titik dari $f_n$ adalah $\frac{1}{1-x}$; yaitu, untuk setiap $x$
dengan $\sabs{x} < h$ bagi suatu $h > 0$,
buktikan bahwa $\lim\limits_{n\to\infty}f_n(x) = \frac{1}{1-x}$.
\end{exercise}

\begin{exercise}
Andaikan $f \colon X \to X$ merupakan kontraksi dengan konstanta kontraksi $k < 1$.  Andaikan Anda menggunakan prosedur
iterasi dengan $x_{n+1} \coloneqq f(x_n)$ seperti dalam pembuktian teorema titik tetap.
Andaikan $x$ merupakan titik
tetap dari $f$.
\begin{enumerate}[a)]
\item
Buktikan bahwa $d(x,x_n) \leq k^n d(x_1,x_0) \frac{1}{1-k}$ untuk semua $n \in \N$.
\item
Andaikan $d(y_1,y_2) \leq 16$ untuk semua $y_1,y_2 \in X$, dan $k=
\nicefrac{1}{2}$.  Tentukan $N$ sedemikian sehingga, dengan memulai dari sebarang titik $x_0 \in X$, 
$d(x,x_n) \leq 2^{-16}$ untuk semua $n \geq N$.
\end{enumerate}
\end{exercise}

\begin{exercise}
Misalkan $f(x) \coloneqq x-\frac{x^2-2}{2x}$ (Anda mungkin mengenal metode Newton untuk
$\sqrt{2}$).
\begin{enumerate}[a)]
\item
Buktikan $f\bigl([1,\infty)\bigr) \subset [1,\infty)$.
\item
Buktikan bahwa $f \colon [1,\infty) \to [1,\infty)$ merupakan kontraksi.
\item
Buktikan bahwa teorema titik tetap dapat diterapkan,
tentukan satu-satunya $x \geq 1$ sedemikian sehingga $f(x) = x$,
dan buktikan bahwa $x = \sqrt{2}$.
Catatan: Secara khusus, teknik dari pembuktian teorema tersebut
dapat digunakan untuk menghampiri $\sqrt{2}$.
\end{enumerate}
\end{exercise}

\begin{exercise}
Andaikan $f \colon X \to X$ merupakan kontraksi dan $(X,d)$ ruang metrik tak kosong
dengan metrik diskret, yaitu $d(x,y) = 1$ untuk setiap $x \neq y$.
Buktikan bahwa $f$ konstan; yaitu,
terdapat $c \in X$ sedemikian sehingga $f(x) = c$ untuk semua $x \in X$.
\end{exercise}

\begin{exercise}
Andaikan $(X,d)$ ruang metrik lengkap tak kosong,
$f \colon X \to X$ suatu pemetaan, dan tuliskan
$f^n$ untuk iterasi ke-$n$ dari $f$.  Andaikan
untuk setiap $n$ terdapat $k_n > 0$ sedemikian sehingga
$d\bigl(f^n(x),f^n(y)\bigr) \leq k_n \, d(x,y)$
untuk semua $x,y \in X$,
dengan $\sum_{n=1}^\infty k_n < \infty$.
Buktikan bahwa $f$ memiliki tepat satu titik tetap dalam $X$.
\end{exercise}
